\documentclass[11pt,oneside]{article}

\usepackage[a4paper,margin=27mm,headheight=15pt]{geometry}
\usepackage[T1]{fontenc}
\usepackage[utf8]{inputenc}
\IfFileExists{libertinus.sty}{\usepackage{libertinus}}{\usepackage{lmodern}}
\usepackage{microtype}
\usepackage{amsmath,amssymb,amsthm,mathtools,mathrsfs}
\usepackage{bm}
\usepackage{booktabs,longtable,array,tabularx}
\usepackage{graphicx}
\usepackage{pgfplots}
\usepackage{xcolor}
\usepackage{enumitem}
\usepackage{fancyhdr}
\usepackage{titlesec}
\usepackage{listings}
\usepackage{xurl}
\usepackage{hyperref}
\usepackage[nameinlink,noabbrev]{cleveref}

\definecolor{linkblue}{RGB}{25,75,135}
\definecolor{shadegray}{RGB}{246,247,249}
\definecolor{rulegray}{RGB}{195,201,208}
\pgfplotsset{compat=1.18}
\hypersetup{
  colorlinks=true,
  linkcolor=linkblue,
  citecolor=linkblue,
  urlcolor=linkblue,
  pdftitle={A Unified Formula Atlas for the Thue--Morse Sequence},
  pdfsubject={Digital, valuative, Pascal, automatic, finite-difference, spectral, and analytic representations},
  pdfkeywords={Thue-Morse, Prouhet, binary digit sum, Lucas theorem, Kummer theorem, Pascal triangle, Mahler equation, Walsh transform, Riesz product, Fabius function}
}

\pagestyle{fancy}
\fancyhf{}
\fancyhead[L]{\small Thue--Morse formula atlas}
\fancyhead[R]{\small\nouppercase{\leftmark}}
\fancyfoot[C]{\thepage}
\renewcommand{\headrulewidth}{0.4pt}

\titleformat{\section}{\Large\bfseries}{\thesection}{0.7em}{}
\titleformat{\subsection}{\large\bfseries}{\thesubsection}{0.7em}{}
\titleformat{\subsubsection}{\normalsize\bfseries}{\thesubsubsection}{0.7em}{}
\setcounter{secnumdepth}{3}
\setcounter{tocdepth}{2}
\setlist[itemize]{leftmargin=2em,itemsep=0.25em,topsep=0.4em}
\setlist[enumerate]{leftmargin=2.3em,itemsep=0.25em,topsep=0.4em}
\setlength{\emergencystretch}{2em}
\allowdisplaybreaks
\numberwithin{equation}{section}

\newtheorem{theorem}{Theorem}[section]
\newtheorem{proposition}[theorem]{Proposition}
\newtheorem{lemma}[theorem]{Lemma}
\newtheorem{corollary}[theorem]{Corollary}
\newtheorem{identity}[theorem]{Identity}
\newtheorem{conjecture}[theorem]{Conjecture}
\theoremstyle{definition}
\newtheorem{definition}[theorem]{Definition}
\newtheorem{algorithm}[theorem]{Algorithm}
\newtheorem{example}[theorem]{Example}
\theoremstyle{remark}
\newtheorem{remark}[theorem]{Remark}
\newtheorem{warning}[theorem]{Warning}

\newcommand{\R}{\mathbb R}
\newcommand{\C}{\mathbb C}
\newcommand{\Q}{\mathbb Q}
\newcommand{\N}{\mathbb N}
\newcommand{\Z}{\mathbb Z}
\newcommand{\Ftwo}{\mathbb F_2}
\newcommand{\e}{\mathrm e}
\newcommand{\ii}{\mathrm i}
\newcommand{\dd}{\,\mathrm d}
\newcommand{\bigO}{\mathcal O}
\newcommand{\smallo}{o}
\newcommand{\sgn}{\operatorname{sgn}}
\newcommand{\sinc}{\operatorname{sinc}}
\newcommand{\wt}{\operatorname{wt}_2}
\newcommand{\vtwo}{\nu_2}
\newcommand{\eps}{\varepsilon}
\newcommand{\xor}{\mathbin{\oplus}}
\newcommand{\submask}{\mathrel{\preccurlyeq_2}}
\newcommand{\ind}{\mathbf 1}
\newcommand{\Bell}{\mathcal B}
\newcommand{\Cat}{\operatorname{Cat}}
\newcommand{\fall}[2]{#1^{\underline{#2}}}
\newcommand{\coeff}[2]{[z^{#1}]\,#2}
\newcommand{\repo}{\nolinkurl{Analysis/FabiusFunction}}
\newcommand{\lean}[1]{%
  \ifmmode\text{\nolinkurl{#1}}\else\nolinkurl{#1}\fi}

\newenvironment{proofidea}{\par\smallskip\noindent\textit{Proof architecture.}\ }{\hfill$\square$\par\smallskip}
\newenvironment{boxedremark}{%
  \par\smallskip\noindent\begin{center}\begin{minipage}{0.94\textwidth}%
  \color{black}\hrule\vspace{0.6em}\small
}{%
  \vspace{0.6em}\hrule\end{minipage}\end{center}\smallskip
}

\lstdefinestyle{pseudo}{
  basicstyle=\ttfamily\small,
  backgroundcolor=\color{shadegray},
  frame=single,
  rulecolor=\color{rulegray},
  columns=fullflexible,
  keepspaces=true,
  showstringspaces=false,
  xleftmargin=0.7em,
  xrightmargin=0.7em,
  aboveskip=0.8em,
  belowskip=0.8em
}

\title{\bfseries A Unified Formula Atlas for the Thue--Morse Sequence\\[0.35em]
\large Digital, valuative, Pascal, automatic, finite-difference, spectral, and analytic representations}
\author{Merged and expanded companion report to the Fabius--Rvachev formal development\\
\small Vladimir Reshetnikov's \texttt{ProveIt} repository}
\date{26 August 2026\\
\small source snapshot at commit \nolinkurl{29e5861c6109207739c41754e1e8bcfce32e84f1}}

\begin{document}
\maketitle

\begin{abstract}
Four draft reports in the Fabius-function research-frontier directory developed
partly overlapping collections of formulae for the Thue--Morse sequence.  This
document merges them into one proof-oriented atlas, removes duplicated
arguments, reconciles notation and sign conventions, and retains the strongest
version of each genuinely different result.  The signed sequence
\(\eps_n=(-1)^{\wt(n)}\) and the zero--one sequence
\(\tau(n)=\wt(n)\bmod2\) are treated simultaneously.

The resulting network begins with binary recursion, block concatenation, XOR
characters, floor and fractional-part sums, trigonometric signs, factorial and
central-binomial valuations, Catalan valuations, and Kummer carry cocycles.  It
then develops the Pascal-row viewpoint in several directions: Glaisher's odd
entry count, the existing exact binomial--logarithm formula, a logarithm-free
modulo-three formula, sparse power-of-two probes, Boolean M\"obius inversion,
multinomial support counts, and a sparse Prouhet theorem on the odd support of
an arbitrary Pascal row.

The lacunary product
\[
  \sum_{n\ge0}\eps_nz^n=\prod_{j\ge0}(1-z^{2^j})
\]
serves as the central analytic engine.  From it we obtain Mahler equations,
cyclotomic factorizations, a ruler-function Euler transform, convolution,
partition, Bell-polynomial, and Hessenberg-determinant formulae, an algebraic
formal series over \(\mathbb F_2\), an integral algebraic lift, exact finite
differences, all dyadic power moments, Walsh and ordinary Fourier transforms,
Riesz products, autocorrelation recurrences, iterated prefix sums, and the
bridge back to the Fabius and Rvachev functions.

New material added in the unification includes an explicit dyadic
weight-distribution chapter, a systematic centered-moment parity theorem, and a
Mellin--Dirichlet theory.  In particular, the shifted Dirichlet series
\[
  D(s,a)=\sum_{n\ge0}\frac{\eps_n}{(n+a)^s}
\]
continues to an entire function, satisfies an exact dyadic parameter equation,
and vanishes at every nonpositive integer.  The proof is driven by the
super-polynomial flatness
\[
 \log\!\prod_{j\ge0}(1-\e^{-2^jt})
 =-\frac{\log^2(1/t)}{2\log2}+\bigO(\log(1/t))
 \qquad(t\downarrow0).
\]
Every elementary assertion is proved in the text; deeper historical results
about transcendence, substitution spectra, and evaluated products are clearly
identified and cited.  No mathematical priority is claimed for identities
merely derived here, and no formula is silently represented as already
formalized in Lean.
\end{abstract}

\tableofcontents
\newpage

\section[Four-draft merge]{Scope, conventions, and the four-draft merge}
\label{sec:scope}

\subsection{Two normalizations}

Write the binary expansion of \(n\in\N=\{0,1,2,\ldots\}\) as
\begin{equation}
  n=\sum_{j\ge0}b_j(n)2^j,
  \qquad b_j(n)\in\{0,1\},
  \label{eq:binary-expansion}
\end{equation}
with only finitely many nonzero digits.  Its binary weight and the two standard
normalizations of the Thue--Morse sequence are
\begin{equation}
  \wt(n)=\sum_{j\ge0}b_j(n),
  \qquad
  \tau(n)=\wt(n)\bmod2\in\{0,1\},
  \qquad
  \eps_n=(-1)^{\wt(n)}=1-2\tau(n).
  \label{eq:normalizations}
\end{equation}
Thus
\begin{equation}
  \tau(n)=\frac{1-\eps_n}{2},
  \qquad
  \eps_n=1-2\tau(n).
  \label{eq:normalization-bridge}
\end{equation}
The signed normalization is usually cleaner for products, transforms, and
cancellation; the zero--one normalization is natural for automata, base
expansions, and the formula in the source exposition.

We use \(a\xor b\) for bitwise exclusive-or and \(k\submask n\) when every
one-bit of \(k\) is also a one-bit of \(n\).  Empty products are \(1\), empty
sums are \(0\), and \(0^0=1\) when it occurs in a finite moment identity.
The valuation \(\vtwo(m)\) is the exponent of \(2\) in a positive integer
\(m\), with \(\vtwo(1)=0\).

\subsection{The Pascal-row formula being complemented}

Section~11.8 of the main Fabius--Rvachev exposition defines
\begin{equation}
  \mathcal X(n)=\sum_{k=0}^{n}(-1)^{\binom nk}
  \label{eq:source-X}
\end{equation}
and proves the exact-power identity
\begin{equation}
  n+1-\mathcal X(n)=2^{\wt(n)+1}.
  \label{eq:source-power}
\end{equation}
Consequently
\begin{equation}
  \tau(n)=
  \frac{1+(-1)^{\log_2\!\left(n+1-
       \sum_{k=0}^{n}(-1)^{\binom nk}\right)}}{2}.
  \label{eq:source-binomial-log}
\end{equation}
The logarithm is exact because its argument was first proved to be a power of
two.  The formula hides the binary expansion inside the parity pattern of a
whole Pascal row.  The present atlas asks what other structures carry the same
parity character.

\subsection{Merge policy}

The four input drafts were not independent papers: large parts shared the same
binary, valuation, generating-product, Prouhet, and Fourier skeleton.  A literal
concatenation would repeat the same proof several times while obscuring the
interesting differences.  The unified text therefore follows three rules.

\begin{enumerate}
\item A theorem common to several drafts appears once, in the strongest useful
      form, with one proof and cross-references to its later specializations.
\item Material unique to a draft is retained when it is mathematically distinct:
      in particular the modulo-three Pascal formula, multinomial logarithms,
      sparse Pascal Prouhet identities, Boolean M\"obius inversion, determinant
      formula, integer algebraic lift, odious/evil enumerators, autocorrelation,
      Dirichlet integrals, and the explicit Fabius-prefix bridge.
\item Statements added during the merge are marked by context rather than by
      novelty rhetoric.  The report claims derivation, not historical priority.
\end{enumerate}

\begin{boxedremark}
\textbf{Proof-status convention.}
The repository already contains formal developments of the binary recurrence,
finite generating product, Prouhet cancellations, iterated prefixes, and the
binomial--logarithm evaluator.  The paper proofs below make many additional
results natural Lean targets, but do not assert that those declarations already
exist.  A dependency-aware formalization map is given in
\cref{sec:formalization}.
\end{boxedremark}

\subsection{A compact formula index}

The table is a navigation device rather than a substitute for hypotheses and
proofs.  Here \(P_m(z)=\sum_{n<2^m}\eps_nz^n\), and
\(E(z)=\sum_{n\ge0}\eps_nz^n\).

\renewcommand{\arraystretch}{1.18}
\begin{longtable}{@{}p{0.25\textwidth}p{0.68\textwidth}@{}}
\toprule
Viewpoint & Representative exact formula \\
\midrule
\endfirsthead
\toprule
Viewpoint & Representative exact formula \\
\midrule
\endhead
Binary recursion & \(\eps_{2n}=\eps_n\), \(\eps_{2n+1}=-\eps_n\). \\
Block concatenation & \(\eps_{2^mq+r}=\eps_q\eps_r\) for \(0\le r<2^m\). \\
XOR character & \(\eps_{a\xor b}=\eps_a\eps_b\). \\
Digit product & \(\eps_n=\prod_j(1-2b_j(n))\). \\
Floors & \(\eps_n=(-1)^{\sum_{j\ge0}\lfloor n/2^j\rfloor}\). \\
Fractional parts & \(\wt(n)=\sum_{j\ge1}\{n/2^j\}\). \\
Factorials & \(\eps_n=(-1)^{n-\vtwo(n!)}\). \\
Central binomial & \(\eps_n=(-1)^{\vtwo\binom{2n}{n}}\). \\
Carries & \(\eps_{a+b}=\eps_a\eps_b(-1)^{\vtwo\binom{a+b}{a}}\). \\
Pascal support & \(\#\{k:\binom nk\text{ odd}\}=2^{\wt(n)}\). \\
Sparse Pascal probes & \(\eps_n=(-1)^{\sum_j\binom{n}{2^j}}\). \\
Log-free Pascal residue & \(\tau(n)=\bigl(2^{\wt(n)}\bmod3\bigr)-1\). \\
Multinomial support & \(O_r(n)=r^{\wt(n)}\). \\
Boolean M\"obius & \(\mu_2(k,n)=\eps_{n\xor k}\) for \(k\submask n\). \\
Finite product & \(P_m(z)=\prod_{j=0}^{m-1}(1-z^{2^j})\). \\
Infinite product & \(E(z)=\prod_{j\ge0}(1-z^{2^j})\), \(|z|<1\). \\
Mahler equation & \(E(z)=(1-z)E(z^2)\). \\
Euler transform & \(\log E(z)=-\sum_{r\ge1}(2^{\vtwo(r)+1}-1)z^r/r\). \\
Finite differences & \(\sum_{n<2^m}\eps_nf(x+n)=(-1)^m\Delta_1\cdots\Delta_{2^{m-1}}f(x)\). \\
Prouhet & \(\sum_{n<2^m}\eps_np(n)=0\) when \(\deg p<m\). \\
Sharp moment & \(\sum_{n<2^m}\eps_nn^m=(-1)^mm!2^{\binom m2}\). \\
Sparse Prouhet & \(\sum_{k\submask n}\eps_kp(k)=0\) when \(\deg p<\wt(n)\). \\
Prefix discrepancy & \(\sum_{n<N}\eps_n=0\) for even \(N\), \(\eps_{(N-1)/2}\) for odd \(N\). \\
Walsh transform & A length-\(2^m\) block is one Walsh frequency. \\
Ordinary DFT & \(\widehat\eps_m(k)=\prod_{j<m}(1-\e^{-2\pi\ii k2^j/2^m})\). \\
Riesz density & \(2^{-m}|P_m(\e^{2\pi\ii x})|^2=\prod_{j<m}(1-\cos(2\pi2^jx))\). \\
Characteristic two & \((1+x)^3T(x)^2+(1+x)^2T(x)+x=0\) in \(\Ftwo[[x]]\). \\
Dirichlet--Mellin & \(D(s,a)=\Gamma(s)^{-1}\int_0^\infty t^{s-1}\e^{-at}E(\e^{-t})\dd t\). \\
\bottomrule
\end{longtable}
\renewcommand{\arraystretch}{1}

\section[Binary character and automata]{Binary character, blocks, automata, and enumerators}
\label{sec:binary}

\subsection{Recursion and uniqueness}

\begin{theorem}[Binary recursion and uniqueness]
\label{thm:binary-recursion}
The signed sequence satisfies
\begin{equation}
  \eps_0=1,
  \qquad
  \eps_{2n}=\eps_n,
  \qquad
  \eps_{2n+1}=-\eps_n.
  \label{eq:eps-recursion}
\end{equation}
Equivalently,
\begin{equation}
  \tau(0)=0,
  \qquad
  \tau(2n)=\tau(n),
  \qquad
  \tau(2n+1)=1-\tau(n).
  \label{eq:tau-recursion}
\end{equation}
Either system and its initial value determines the sequence uniquely.
\end{theorem}

\begin{proof}
Appending a terminal binary digit \(0\) leaves the digit sum unchanged;
appending a terminal \(1\) increases it by one.  For uniqueness, every positive
integer is \(2q\) or \(2q+1\) with \(q<n\), so strong induction eventually
reduces its value to the initial value at zero.
\end{proof}

Let \(W_m=\tau(0)\tau(1)\cdots\tau(2^m-1)\), and let
\(\overline W\) denote bitwise complement.  The recurrence gives
\begin{equation}
  W_{m+1}=W_m\,\overline{W_m}.
  \label{eq:word-recursion}
\end{equation}
Thus the infinite word is the fixed point beginning with \(0\) of
\begin{equation}
  0\longmapsto01,
  \qquad
  1\longmapsto10.
  \label{eq:thue-morphism}
\end{equation}
In signed form a block is followed by its negative.

\subsection{Binary concatenation and reflection}

\begin{theorem}[Block concatenation]
\label{thm:block-concatenation}
For \(m,q\ge0\) and \(0\le r<2^m\),
\begin{align}
  \wt(2^mq+r)&=\wt(q)+\wt(r),
  \label{eq:block-weight}\\
  \eps_{2^mq+r}&=\eps_q\eps_r,
  \label{eq:block-sign}\\
  \tau(2^mq+r)&=\tau(q)\xor\tau(r).
  \label{eq:block-bit}
\end{align}
\end{theorem}

\begin{proof}
Multiplication by \(2^m\) shifts the binary expansion of \(q\) by \(m\)
places.  Since \(r<2^m\), adding \(r\) fills only those vacated positions and
creates no carry.  The bits concatenate, so weights add and parities follow.
\end{proof}

\begin{corollary}[Reflection and bit-position symmetry]
\label{cor:block-reflection}
For \(0\le n<2^m\),
\begin{equation}
  \eps_{2^m-1-n}=(-1)^m\eps_n,
  \qquad
  \tau(2^m-1-n)=\tau(n)\xor(m\bmod2).
  \label{eq:block-reflection}
\end{equation}
Permuting the first \(m\) binary positions leaves \(\eps_n\) and \(\tau(n)\)
unchanged.
\end{corollary}

\begin{proof}
Within \(m\) positions, \(2^m-1-n\) is the bitwise complement of \(n\), so
its weight is \(m-\wt(n)\).  A permutation of positions leaves their sum
unchanged.
\end{proof}

\subsection{Exclusive-or is the natural group law}

\begin{theorem}[XOR character]
\label{thm:xor-character}
For all \(a,b\ge0\),
\begin{equation}
  \tau(a\xor b)=\tau(a)\xor\tau(b),
  \qquad
  \boxed{\eps_{a\xor b}=\eps_a\eps_b.}
  \label{eq:xor-character}
\end{equation}
Thus \(n\mapsto\eps_n\) is a character of the Boolean group
\(\bigoplus_{j\ge0}\Z/2\Z\).
\end{theorem}

\begin{proof}
At each position, XOR adds the two bits modulo two.  Summing the coordinate
identities modulo two gives
\(\wt(a\xor b)\equiv\wt(a)+\wt(b)\pmod2\).
\end{proof}

\begin{corollary}[XOR autocorrelation and Hamming spheres]
\label{cor:xor-hamming}
Fix \(m\ge0\), \(0\le a<2^m\), and \(0\le r\le m\).  Then
\begin{align}
  \sum_{n=0}^{2^m-1}\eps_n\eps_{n\xor a}
   &=2^m\eps_a,
  \label{eq:xor-autocorrelation}\\
  \sum_{\substack{0\le n<2^m\\\wt(n\xor a)=r}}\eps_n
   &=\eps_a(-1)^r\binom mr.
  \label{eq:hamming-sphere}
\end{align}
\end{corollary}

\begin{proof}
The first summand is termwise \(\eps_a\).  For the second, put
\(h=n\xor a\).  Then \(\eps_n=\eps_a\eps_h=\eps_a(-1)^r\), and there are
\(\binom mr\) masks of weight \(r\).
\end{proof}

\subsection{Finite automata, matrices, and tensor powers}

The recurrence gives a two-state automaton: begin in state \(+1\), read the
binary digits in either order, and flip the state for each one-bit.  In matrix
form, let
\begin{equation}
  M_0=\begin{pmatrix}1&0\\0&1\end{pmatrix},
  \qquad
  M_1=\begin{pmatrix}0&1\\1&0\end{pmatrix},
  \qquad
  e_0=\binom10.
  \label{eq:automaton-matrices}
\end{equation}
If \(n=\sum_{j=0}^{m-1}b_j2^j\), then
\begin{equation}
  M_{b_0}\cdots M_{b_{m-1}}e_0
  =\binom{1-\tau(n)}{\tau(n)}.
  \label{eq:automaton-output}
\end{equation}
The matrices commute, reflecting the fact that only the number of one-bits
matters.  In natural binary order,
\begin{equation}
  (\eps_0,\ldots,\eps_{2^m-1})^{\mathsf T}
  =\bigl((1,-1)^{\mathsf T}\bigr)^{\otimes m}.
  \label{eq:tensor-block}
\end{equation}
Hence every residue subsequence in the \(2\)-kernel is either \(\eps\) or
\(-\eps\):
\begin{equation}
  (\eps_{2^mn+r})_{n\ge0}=\eps_r(\eps_n)_{n\ge0}.
  \label{eq:two-kernel}
\end{equation}
This is the finite-kernel criterion for \(2\)-automaticity in its smallest
possible form.

\begin{algorithm}[XOR-fold evaluator]
\label{alg:xor-fold}
For a fixed-width machine word, fold the bits by XORing with successively
larger right shifts.  The final least significant bit is \(\tau(n)\).
\begin{lstlisting}[style=pseudo]
x := n
x := x xor (x >> 1)
x := x xor (x >> 2)
x := x xor (x >> 4)
x := x xor (x >> 8)
... continue while the shift is below the word width ...
tau := x & 1
epsilon := 1 - 2*tau
\end{lstlisting}
\end{algorithm}

\subsection{Exact weight distribution in a dyadic block}
\label{sec:weight-distribution}

The four drafts repeatedly use the fact that a length-\(2^m\) block is a
Boolean cube.  It is useful to record the complete distribution, not only its
parity consequence.

\begin{theorem}[Binomial weight enumerator]
\label{thm:weight-enumerator}
For every \(m\ge0\),
\begin{equation}
  \boxed{
  \sum_{n=0}^{2^m-1}u^{\wt(n)}=(1+u)^m.}
  \label{eq:weight-enumerator}
\end{equation}
Consequently
\begin{equation}
  \#\{0\le n<2^m:\wt(n)=r\}=\binom mr.
  \label{eq:weight-distribution}
\end{equation}
Under the uniform distribution on the block, \(\wt(n)\) is binomial with
mean \(m/2\) and variance \(m/4\).
\end{theorem}

\begin{proof}
Each of the \(m\) positions independently contributes \(1\) or \(u\).
Multiplication gives \((1+u)^m\); coefficient extraction and the standard
binomial moments give the remaining statements.
\end{proof}

Setting \(u=-1\) gives exact sign balance for \(m\ge1\).  More generally, the
characteristic function of the weight is
\begin{equation}
  2^{-m}\sum_{n<2^m}\e^{\ii t\wt(n)}
  =\left(\frac{1+\e^{\ii t}}2\right)^m.
  \label{eq:weight-characteristic}
\end{equation}
This elementary statistical layer is useful when comparing Boolean/Walsh
structure with ordinary Fourier structure later in the report.

\subsection{Evil and odious enumerators}

Call \(n\) \emph{evil} when \(\wt(n)\) is even and \emph{odious} when it is
odd.  Let \(\mathcal E_n\) and \(\mathcal O_n\) be the increasing enumerations,
starting at index zero.

\begin{theorem}[Evil and odious numbers]
\label{thm:evil-odious}
For every \(n\ge0\),
\begin{equation}
  \boxed{
  \mathcal E_n=2n+\tau(n),
  \qquad
  \mathcal O_n=2n+1-\tau(n).}
  \label{eq:evil-odious}
\end{equation}
Hence
\begin{equation}
  \mathcal O_n-\mathcal E_n=\eps_n,
  \qquad
  \mathcal O_n+\mathcal E_n=4n+1.
  \label{eq:evil-odious-sumdiff}
\end{equation}
\end{theorem}

\begin{proof}
The two consecutive integers \(2n\) and \(2n+1\) have parities
\(\tau(n)\) and \(1-\tau(n)\) of their binary weights.  If \(\tau(n)=0\),
then \(2n\) is evil and \(2n+1\) odious; if \(\tau(n)=1\), the roles reverse.
The displayed formulas choose the correct member of each consecutive pair.
\end{proof}

\begin{corollary}[Composition laws]
\label{cor:evil-odious-composition}
For every \(n\ge0\),
\begin{align}
  \mathcal E_{\mathcal E_n}&=2\mathcal E_n,
  &\mathcal O_{\mathcal E_n}&=2\mathcal E_n+1,
  \label{eq:EE-OE}\\
  \mathcal E_{\mathcal O_n}&=2\mathcal O_n+1,
  &\mathcal O_{\mathcal O_n}&=2\mathcal O_n.
  \label{eq:EO-OO}
\end{align}
\end{corollary}

\begin{proof}
By definition \(\tau(\mathcal E_n)=0\) and
\(\tau(\mathcal O_n)=1\).  Substitute these values into
\eqref{eq:evil-odious} with the new index.
\end{proof}

\section[Digit and trigonometric formulae]{Digit, floor, fractional-part, and trigonometric formulas}
\label{sec:digits}

\subsection{Extracting one binary digit}

\begin{lemma}[Floor extraction]
\label{lem:bit-floor}
For every \(j,n\ge0\),
\begin{equation}
  b_j(n)=\left\lfloor\frac{n}{2^j}\right\rfloor
          -2\left\lfloor\frac{n}{2^{j+1}}\right\rfloor.
  \label{eq:bit-floor}
\end{equation}
\end{lemma}

\begin{proof}
Write \(\lfloor n/2^j\rfloor=2q+r\) with \(r\in\{0,1\}\).  Then
\(q=\lfloor n/2^{j+1}\rfloor\), while \(r\) is exactly the \(j\)-th binary
digit.
\end{proof}

Since \((-1)^b=1-2b\) for \(b\in\{0,1\}\), the definition factorizes over
coordinates.

\begin{theorem}[Digit product and Boolean polynomial]
\label{thm:digit-product}
For every \(n\ge0\),
\begin{equation}
  \boxed{
  \eps_n=\prod_{j\ge0}\bigl(1-2b_j(n)\bigr),
  \qquad
  \tau(n)=\frac{1-\prod_{j\ge0}(1-2b_j(n))}{2}.}
  \label{eq:digit-product}
\end{equation}
Only finitely many factors differ from \(1\).  If \(J\) is any finite set
containing the support of the binary expansion, then
\begin{equation}
  \boxed{
  \tau(n)=\sum_{\varnothing\ne S\subseteq J}
     (-2)^{|S|-1}\prod_{j\in S}b_j(n).}
  \label{eq:boolean-polynomial}
\end{equation}
\end{theorem}

\begin{proof}
Multiply \((-1)^{b_j}=1-2b_j\) over all positions.  Expanding the product in
the zero--one conversion \eqref{eq:normalization-bridge} gives
\eqref{eq:boolean-polynomial}.
\end{proof}

Substitution of \eqref{eq:bit-floor} gives a digit-free, floor-only product:
\begin{equation}
  \eps_n=
  \prod_{j\ge0}
  \left(
    1-2\left\lfloor\frac{n}{2^j}\right\rfloor
      +4\left\lfloor\frac{n}{2^{j+1}}\right\rfloor
  \right).
  \label{eq:pure-floor-product}
\end{equation}
Every factor is \(1-2b_j(n)\in\{1,-1\}\).

\subsection{Dyadic floor and fractional-part sums}

\begin{lemma}[Legendre floor sum]
\label{lem:floor-sum}
For every \(n\ge0\),
\begin{align}
  \sum_{j\ge1}\left\lfloor\frac{n}{2^j}\right\rfloor
   &=n-\wt(n),
  \label{eq:floor-sum}\\
  \sum_{j\ge0}\left\lfloor\frac{n}{2^j}\right\rfloor
   &=2n-\wt(n).
  \label{eq:floor-sum-all}
\end{align}
\end{lemma}

\begin{proof}
Insert the binary expansion and interchange finite sums:
\begin{align*}
 \sum_{j\ge1}\left\lfloor\frac n{2^j}\right\rfloor
 &=\sum_{j\ge1}\sum_{\ell\ge j}b_\ell(n)2^{\ell-j}\\
 &=\sum_{\ell\ge1}b_\ell(n)(2^\ell-1)
  =n-\sum_{\ell\ge0}b_\ell(n).
\end{align*}
Adding the \(j=0\) term gives the second identity.
\end{proof}

\begin{theorem}[Floor-sum formulae]
\label{thm:floor-formulas}
For every \(n\ge0\),
\begin{equation}
  \boxed{
  \eps_n=(-1)^{\sum_{j\ge0}\lfloor n/2^j\rfloor},
  \qquad
  \tau(n)\equiv
     \sum_{j\ge0}\left\lfloor\frac{n}{2^j}\right\rfloor
     \pmod2.}
  \label{eq:floor-formulas}
\end{equation}
Equivalently,
\begin{equation}
  \eps_n=\prod_{j\ge0}(-1)^{\lfloor n/2^j\rfloor}.
  \label{eq:floor-sign-product}
\end{equation}
\end{theorem}

\begin{proof}
By \eqref{eq:floor-sum-all}, the exponent differs from \(\wt(n)\) by an even
integer.
\end{proof}

There is also a genuinely infinite real-valued representation whose sum is
nevertheless always an integer.

\begin{proposition}[Fractional-part digit sum]
\label{prop:fractional-weight}
For every \(n\ge0\),
\begin{equation}
  \boxed{
  \wt(n)=\sum_{j\ge1}\left\{\frac{n}{2^j}\right\}.}
  \label{eq:fractional-weight}
\end{equation}
Consequently
\begin{equation}
  \eps_n=
  \cos\!\left(
    \pi\sum_{j\ge1}\left\{\frac n{2^j}\right\}
  \right),
  \qquad
  \tau(n)=\frac{1-\eps_n}{2}.
  \label{eq:fractional-cos}
\end{equation}
\end{proposition}

\begin{proof}
The geometric series gives \(\sum_{j\ge1}n/2^j=n\).  Subtract
\eqref{eq:floor-sum} termwise:
\[
 \sum_{j\ge1}\left\{\frac n{2^j}\right\}
 =n-\sum_{j\ge1}\left\lfloor\frac n{2^j}\right\rfloor
 =\wt(n).
\]
\end{proof}

\subsection{Square waves and a sine-sign product}

For nonintegral \(x\),
\begin{equation}
  \sgn(\sin\pi x)=(-1)^{\lfloor x\rfloor}.
  \label{eq:square-wave}
\end{equation}
The half-unit shift avoids every zero.

\begin{theorem}[Rademacher--sine formula]
\label{thm:sine-sign}
For every \(n\ge0\),
\begin{equation}
  \boxed{
  \eps_n=\prod_{j\ge0}
  \sgn\!\left(
    \sin\frac{\pi(n+\tfrac12)}{2^j}
  \right).}
  \label{eq:sine-sign}
\end{equation}
All sufficiently large factors are \(+1\).  Moreover the \(j\)-th factor is
exactly \((-1)^{b_j(n)}\).
\end{theorem}

\begin{proof}
The number \((n+\tfrac12)/2^j\) is never integral and has the same floor as
\(n/2^j\).  Apply \eqref{eq:square-wave}, then observe from
\eqref{eq:bit-floor} that
\((-1)^{\lfloor n/2^j\rfloor}=(-1)^{b_j(n)}\).  The factors are positive once
\(2^j>n+\tfrac12\).
\end{proof}

\subsection{The sign of the Fabius sinc product}
\label{sec:sinc-sign}

Define \(\sinc y=\sin y/y\), with \(\sinc0=1\), and
\begin{equation}
  \Phi(x)=\prod_{j=0}^{\infty}
    \sinc\!\left(\frac{\pi x}{2^j}\right).
  \label{eq:sinc-product}
\end{equation}
The product converges locally uniformly because the logarithm of the tail is
controlled by \(\sum_jx^2/4^j\).  It is, up to normalization conventions, the
infinite Fourier product naturally associated with the Rvachev up-function.

\begin{corollary}[Interval sign of the sinc product]
\label{cor:sinc-sign}
If \(N\in\N\) and \(N<x<N+1\), then
\begin{equation}
  \sgn\Phi(x)=\eps_N.
  \label{eq:sinc-interval-sign}
\end{equation}
In particular,
\begin{equation}
  \eps_n=\sgn\Phi\!\left(n+\tfrac12\right).
  \label{eq:sinc-half-sample}
\end{equation}
\end{corollary}

\begin{proof}
For \(x>0\), every sinc denominator is positive.  No multiple of \(2^j\) lies
strictly between the consecutive integers \(N\) and \(N+1\), so
\(\lfloor x/2^j\rfloor=\lfloor N/2^j\rfloor\).  The numerator signs therefore
multiply to
\((-1)^{\sum_j\lfloor N/2^j\rfloor}=\eps_N\) by
\cref{thm:floor-formulas}.
\end{proof}

\subsection{The ruler transition}

Incrementing \(n\) flips all trailing one-bits to zero and the next zero-bit
to one.  The number of erased one-bits is \(\vtwo(n+1)\).

\begin{theorem}[Successor and ruler formulae]
\label{thm:ruler-transition}
For every \(n\ge0\),
\begin{align}
  \wt(n+1)-\wt(n)&=1-\vtwo(n+1),
  \label{eq:weight-successor}\\
  \eps_{n+1}&=(-1)^{1+\vtwo(n+1)}\eps_n,
  \label{eq:eps-successor}\\
  (-1)^{\vtwo(n+1)}&=-\eps_n\eps_{n+1}.
  \label{eq:ruler-adjacent}
\end{align}
Consequently
\begin{equation}
  \boxed{
  \eps_n=\prod_{k=1}^{n}(-1)^{1+\vtwo(k)}.}
  \label{eq:ruler-product}
\end{equation}
\end{theorem}

\begin{proof}
The bit transition changes the weight by one minus the number of trailing
ones.  Exponentiate \((-1)\), then telescope from \(\eps_0=1\).
\end{proof}

\section[Valuations and carries]{Factorial valuations, central binomials, Catalan numbers, and carries}
\label{sec:valuations}

\subsection{Legendre's formula}

The floor sum in \cref{lem:floor-sum} counts powers of two in a factorial.

\begin{theorem}[Factorial valuation]
\label{thm:factorial-valuation}
For every \(n\ge0\),
\begin{equation}
  \vtwo(n!)=\sum_{j\ge1}\left\lfloor\frac n{2^j}\right\rfloor
           =n-\wt(n).
  \label{eq:legendre-two}
\end{equation}
Hence
\begin{equation}
  \boxed{
  \eps_n=(-1)^{n-\vtwo(n!)}=(-1)^{n+\vtwo(n!)},
  \qquad
  \tau(n)\equiv n+\vtwo(n!)\pmod2.}
  \label{eq:factorial-thue-morse}
\end{equation}
\end{theorem}

\begin{proof}
Every multiple of \(2\) contributes at least one factor, every multiple of
\(4\) contributes one additional factor, and so on.  This gives the floor
sum; \cref{lem:floor-sum} gives the digit form.  Addition and subtraction agree
modulo two.
\end{proof}

The telescoping product \eqref{eq:ruler-product} returns to this formula because
\(\sum_{k=1}^{n}\vtwo(k)=\vtwo(n!)\).

\subsection{A single central binomial coefficient}

\begin{theorem}[Central-binomial valuation]
\label{thm:central-binomial}
For every \(n\ge0\),
\begin{equation}
  \boxed{
  \vtwo\binom{2n}{n}=\wt(n).}
  \label{eq:central-v2}
\end{equation}
Consequently
\begin{equation}
  \boxed{
  \eps_n=(-1)^{\vtwo\binom{2n}{n}},
  \qquad
  \tau(n)=\vtwo\binom{2n}{n}\bmod2.}
  \label{eq:central-thue-morse}
\end{equation}
Equivalently,
\begin{equation}
  \tau(n)=\sin^2\!\left(\frac\pi2\vtwo\binom{2n}{n}\right),
  \qquad
  \eps_n=\cos\!\left(\pi\vtwo\binom{2n}{n}\right).
  \label{eq:central-trig}
\end{equation}
\end{theorem}

\begin{proof}
Use \eqref{eq:legendre-two} and \(\wt(2n)=\wt(n)\):
\begin{align*}
 \vtwo\binom{2n}{n}
 &=\vtwo((2n)!)-2\vtwo(n!)\\
 &=\bigl(2n-\wt(2n)\bigr)-2\bigl(n-\wt(n)\bigr)
  =\wt(n).
\end{align*}
\end{proof}

\begin{boxedremark}
\textbf{The shortest Pascal replacement.}
The source formula inspects all \(n+1\) entries of row \(n\) and takes an exact
logarithm.  Formula \eqref{eq:central-thue-morse} inspects the two-adic order of
one entry in row \(2n\), and in fact recovers the entire digit sum rather than
only its parity.
\end{boxedremark}

\subsection{Catalan valuations}

Let
\begin{equation}
  \Cat_n=\frac1{n+1}\binom{2n}{n}.
  \label{eq:catalan-definition}
\end{equation}

\begin{corollary}[Catalan formula]
\label{cor:catalan-valuation}
For every \(n\ge0\),
\begin{equation}
  \vtwo(\Cat_n)=\wt(n+1)-1,
  \label{eq:catalan-v2}
\end{equation}
and hence
\begin{equation}
  \boxed{
  \eps_{n+1}=-(-1)^{\vtwo(\Cat_n)}.}
  \label{eq:catalan-thue-morse}
\end{equation}
\end{corollary}

\begin{proof}
By \cref{thm:central-binomial},
\(\vtwo(\Cat_n)=\wt(n)-\vtwo(n+1)\).  Equation
\eqref{eq:weight-successor} rewrites this as \(\wt(n+1)-1\).
\end{proof}

\subsection{Kummer's carry cocycle}

\begin{theorem}[Addition with carries]
\label{thm:carry-cocycle}
For all \(a,b\ge0\),
\begin{equation}
  \vtwo\binom{a+b}{a}
  =\wt(a)+\wt(b)-\wt(a+b).
  \label{eq:kummer-digit-sum}
\end{equation}
This is the number of carry operations in the binary addition of \(a\) and
\(b\), including propagated carries.  Therefore
\begin{equation}
  \boxed{
  \eps_{a+b}=\eps_a\eps_b
      (-1)^{\vtwo\binom{a+b}{a}}.}
  \label{eq:carry-sign}
\end{equation}
Equivalently,
\begin{equation}
  \tau(a+b)\equiv\tau(a)+\tau(b)+
    \vtwo\binom{a+b}{a}\pmod2.
  \label{eq:carry-bit}
\end{equation}
\end{theorem}

\begin{proof}
Legendre's formula gives
\begin{align*}
 \vtwo\binom{a+b}{a}
 &=(a+b-\wt(a+b))-(a-\wt(a))-(b-\wt(b)),
\end{align*}
which is \eqref{eq:kummer-digit-sum}.  Rearranging modulo two and raising
\((-1)\) to both sides gives the sign formula.  The carry interpretation is
the binary case of Kummer's theorem.
\end{proof}

The no-carry case \(a\mathbin{\&}b=0\) makes the valuation zero and recovers
ordinary multiplicativity.  In particular, \cref{thm:block-concatenation} is
the separated-support case.

\begin{corollary}[Multinomial carry formula]
\label{cor:multinomial-carry}
For \(a_1,\ldots,a_r\ge0\) and \(N=a_1+\cdots+a_r\),
\begin{align}
 \vtwo\binom{N}{a_1,\ldots,a_r}
 &=\sum_{j=1}^{r}\wt(a_j)-\wt(N),
 \label{eq:multinomial-v2}\\
 \eps_N
 &=\left(\prod_{j=1}^{r}\eps_{a_j}\right)
   (-1)^{\vtwo\binom{N}{a_1,\ldots,a_r}}.
 \label{eq:multinomial-carry-sign}
\end{align}
\end{corollary}

\begin{proof}
Apply Legendre's formula to
\(N!/(a_1!\cdots a_r!)\) and then take parity.
\end{proof}

\section[Pascal support and Boolean inversion]{Pascal support, residue formulae, multinomials, and Boolean inversion}
\label{sec:pascal}

\subsection{Lucas's theorem as a submask criterion}

\begin{theorem}[Lucas parity criterion]
\label{thm:lucas-submask}
For \(0\le k\le n\),
\begin{equation}
  \binom nk\equiv1\pmod2
  \quad\Longleftrightarrow\quad
  k\submask n.
  \label{eq:lucas-submask}
\end{equation}
Equivalently,
\begin{equation}
  \ind_{k\submask n}
   =\frac{1-(-1)^{\binom nk}}2.
  \label{eq:odd-indicator}
\end{equation}
\end{theorem}

\begin{proof}
In \(\Ftwo[z]\), Frobenius gives
\[
 (1+z)^n
 =\prod_{j:b_j(n)=1}(1+z)^{2^j}
 =\prod_{j:b_j(n)=1}(1+z^{2^j}).
\]
The exponents occurring in the expansion are exactly the submasks of \(n\),
and their coefficients are one modulo two.
\end{proof}

Define the odd-support polynomial of row \(n\) by
\begin{equation}
  Q_n(z)=\sum_{k=0}^{n}
       \ind_{\binom nk\text{ odd}}z^k.
  \label{eq:row-support-polynomial}
\end{equation}
The proof gives the complete factorization
\begin{equation}
  \boxed{
  Q_n(z)=\prod_{j:b_j(n)=1}(1+z^{2^j}).}
  \label{eq:row-support-factorization}
\end{equation}

\begin{corollary}[Glaisher--Fine count]
\label{cor:glaisher-count}
The number \(\nu(n)\) of odd entries in row \(n\) is
\begin{equation}
  \boxed{\nu(n)=2^{\wt(n)}.}
  \label{eq:glaisher-count}
\end{equation}
\end{corollary}

\begin{proof}
Set \(z=1\) in \eqref{eq:row-support-factorization}.
\end{proof}

The source formula is now transparent.  Since even entries contribute \(+1\)
and odd entries contribute \(-1\),
\begin{equation}
  \mathcal X(n)=n+1-2\nu(n)=n+1-2^{\wt(n)+1},
  \label{eq:X-decomposition}
\end{equation}
which is exactly \eqref{eq:source-power}.

\subsection{A logarithm-free Pascal-row formula modulo three}
\label{sec:mod-three}

Let \(\operatorname{rem}_3(x)\in\{0,1,2\}\) denote the least nonnegative
residue modulo three.  Powers of two alternate between one and two modulo
three.

\begin{theorem}[Odd-count residue formula]
\label{thm:mod-three}
For every \(n\ge0\),
\begin{equation}
  \boxed{
  \tau(n)=\operatorname{rem}_3(\nu(n))-1,
  \qquad
  \eps_n=3-2\operatorname{rem}_3(\nu(n)).}
  \label{eq:odd-count-mod3}
\end{equation}
The residue is always \(1\) or \(2\), never \(0\).
\end{theorem}

\begin{proof}
By \eqref{eq:glaisher-count}, \(\nu(n)=2^{\wt(n)}\).  Its residue is one when
\(\wt(n)\) is even and two when it is odd.
\end{proof}

A form using exactly the row signs of the source formula can be written without
the odd-count notation.  Put
\begin{equation}
  S(n)=\sum_{k=1}^{n}(-1)^{\binom nk},
  \qquad S(0)=0.
  \label{eq:S-pascal}
\end{equation}
Since the omitted \(k=0\) term is \(-1\),
\begin{equation}
  S(n)=n+2-2^{\wt(n)+1}.
  \label{eq:S-pascal-exact}
\end{equation}

\begin{corollary}[Binomial residue and sine formulae]
\label{cor:mod-three-sine}
For every \(n\ge0\),
\begin{align}
  \boxed{
  \tau(n)=\operatorname{rem}_3\bigl(2n+S(n)\bigr),}
  \label{eq:pascal-mod3}\\
  \boxed{
  \tau(n)=\frac43\sin^2\!\left(
     \frac\pi3\bigl(n-S(n)\bigr)
  \right).}
  \label{eq:pascal-sine}
\end{align}
\end{corollary}

\begin{proof}
Equation \eqref{eq:S-pascal-exact} gives
\(2n+S(n)=3n+2-2^{\wt(n)+1}\), whose residue is zero or one according as the
weight is even or odd.  Also
\(n-S(n)=2^{\wt(n)+1}-2\), congruent to zero or two modulo three; the squared
sine is therefore zero or \(3/4\).
\end{proof}

The logarithmic and residue formulas are two shadows of the same invariant
\(\nu(n)\): the first reads the exponent of a power of two; the second reads
that exponent's parity from the two-cycle of powers modulo three.

\subsection{A one-parameter row statistic}

For a scalar \(u\), define
\begin{equation}
  R_u(n)=\sum_{k=0}^{n}u^{\binom nk\bmod2}.
  \label{eq:Ru-definition}
\end{equation}
Exactly \(\nu(n)\) exponents are one and the others are zero, so
\begin{equation}
  \boxed{
  R_u(n)=n+1+(u-1)2^{\wt(n)}.}
  \label{eq:Ru-exact}
\end{equation}
Thus \(R_{-1}(n)=\mathcal X(n)\), while \(R_0(n)\) counts the even entries.
For every fixed \(u\ne1\),
\begin{equation}
  2^{\wt(n)}=\frac{R_u(n)-(n+1)}{u-1}.
  \label{eq:Ru-recovery}
\end{equation}
The signed choice \(u=-1\) is elegant, but structurally it is one member of a
whole exact family.

\subsection{Sparse Pascal probes for individual bits}

Lucas's criterion at \(k=2^j\) isolates one binary coordinate:
\begin{equation}
  \binom{n}{2^j}\equiv b_j(n)\pmod2.
  \label{eq:pascal-bit-probe}
\end{equation}

\begin{theorem}[Power-of-two column formula]
\label{thm:sparse-pascal-probes}
For every \(n\ge0\),
\begin{align}
  \boxed{
  \eps_n=(-1)^{\sum_{j\ge0}\binom{n}{2^j}}}
  &=\prod_{j\ge0}\cos\!\left(\pi\binom{n}{2^j}\right),
  \label{eq:sparse-pascal-sign}\\
  \tau(n)&=\frac12\left(
      1-\prod_{j\ge0}\cos\!\left(\pi\binom{n}{2^j}\right)
    \right).
  \label{eq:sparse-pascal-bit}
\end{align}
Only the columns \(2^j\le n\) matter.
\end{theorem}

\begin{proof}
Only the parity of each binomial coefficient contributes to the exponent or
cosine.  Equation \eqref{eq:pascal-bit-probe} turns their parity sum into
\(\wt(n)\).
\end{proof}

Unlike the source formula, this scans only \(O(\log n)\) selected entries and
uses no logarithm.  Its price is the distinguished set of power-of-two
columns.

\subsection{Odd multinomial coefficients and exact \texorpdfstring{\(r\)}{r}-nomial logarithms}

Fix \(r\ge2\).  Let \(O_r(n)\) be the number of odd coefficients in
\((x_1+\cdots+x_r)^n\).

\begin{theorem}[Multinomial support count]
\label{thm:multinomial-support}
For every \(n\ge0\),
\begin{equation}
  \boxed{O_r(n)=r^{\wt(n)}.}
  \label{eq:multinomial-support}
\end{equation}
Consequently
\begin{equation}
  \boxed{
  \eps_n=(-1)^{\log_rO_r(n)},
  \qquad
  \tau(n)=\frac{1-(-1)^{\log_rO_r(n)}}2.}
  \label{eq:multinomial-log}
\end{equation}
The logarithm is exact.
\end{theorem}

\begin{proof}
In characteristic two,
\begin{equation}
 (x_1+\cdots+x_r)^n
 =\prod_{j:b_j(n)=1}
   (x_1^{2^j}+\cdots+x_r^{2^j}).
  \label{eq:multinomial-frobenius}
\end{equation}
For each one-bit of \(n\), choose independently one of the \(r\) variables to
receive it.  Unique binary expansions make distinct assignments produce
distinct exponent vectors, so there are \(r^{\wt(n)}\) nonzero coefficients
modulo two.
\end{proof}

There are \(\binom{n+r-1}{r-1}\) weak compositions of \(n\) into \(r\) parts.
Define
\begin{equation}
  \mathcal X_r(n)=
  \sum_{k_1+\cdots+k_r=n}
  (-1)^{\binom{n}{k_1,\ldots,k_r}}.
  \label{eq:Xr-definition}
\end{equation}

\begin{corollary}[Signed \(r\)-nomial row formula]
\label{cor:rnomial-log}
For every \(r\ge2\),
\begin{equation}
  \binom{n+r-1}{r-1}-\mathcal X_r(n)
   =2r^{\wt(n)},
  \label{eq:rnomial-exact}
\end{equation}
and therefore
\begin{equation}
  \boxed{
  \tau(n)=\frac{1-(-1)^{\displaystyle
    \log_r\!\left(
      \frac{\binom{n+r-1}{r-1}-\mathcal X_r(n)}2
    \right)}}2.}
  \label{eq:rnomial-log-formula}
\end{equation}
\end{corollary}

\begin{proof}
Each odd multinomial coefficient contributes \(-1\), each even coefficient
\(+1\).  Subtracting the signed sum from the total count gives twice the odd
support size, which is \(2r^{\wt(n)}\) by
\cref{thm:multinomial-support}.
\end{proof}

\subsection{Boolean weight enumerators and M\"obius inversion}

The submasks of \(n\) form a Boolean lattice on its set of one-bit positions.

\begin{theorem}[Submask weight enumerator]
\label{thm:submask-enumerator}
For every indeterminate \(z\),
\begin{align}
  \boxed{
  \sum_{k\submask n}z^{\wt(k)}=(1+z)^{\wt(n)},}
  \label{eq:submask-enumerator}\\
  \boxed{
  z^{\wt(n)}=\sum_{k\submask n}(z-1)^{\wt(k)}.}
  \label{eq:submask-inversion-polynomial}
\end{align}
More generally,
\begin{equation}
  \sum_{k\submask n}
    x^{\wt(k)}y^{\wt(n-k)}=(x+y)^{\wt(n)}.
  \label{eq:submask-bivariate}
\end{equation}
\end{theorem}

\begin{proof}
For each occupied bit of \(n\), choose whether it belongs to \(k\), to its
complement \(n-k\), or is omitted in the one-variable specialization.
The second formula is the first with \(z\) replaced by \(z-1\).
\end{proof}

Important specializations are
\begin{align}
  \#\{k:k\submask n\}&=2^{\wt(n)},
  \label{eq:submask-count}\\
  \sum_{k\submask n}\eps_k
  &=\begin{cases}1,&n=0,\\0,&n>0,\end{cases}
  \label{eq:submask-cancellation}\\
  \boxed{
  \eps_n=\sum_{k\submask n}(-2)^{\wt(k)}.}
  \label{eq:submask-reconstruction}
\end{align}

\begin{theorem}[Boolean M\"obius inversion]
\label{thm:boolean-mobius}
If \(k\submask n\), then the M\"obius function of the submask poset is
\begin{equation}
  \mu_2(k,n)=(-1)^{\wt(n)-\wt(k)}=\eps_{n\xor k}.
  \label{eq:boolean-mobius}
\end{equation}
Thus
\begin{equation}
  g(n)=\sum_{k\submask n}f(k)
  \quad\Longleftrightarrow\quad
  \boxed{
  f(n)=\sum_{k\submask n}\eps_{n\xor k}g(k).}
  \label{eq:boolean-mobius-inversion}
\end{equation}
\end{theorem}

\begin{proof}
The interval \([k,n]\) is the Boolean lattice on the bits of \(n\) absent
from \(k\); its rank is \(\wt(n)-\wt(k)=\wt(n\xor k)\).  The usual Boolean
M\"obius function is \((-1)^{\text{rank}}\), and standard M\"obius inversion
applies.
\end{proof}

Let
\begin{equation}
  P_{n,k}=\ind_{\binom nk\text{ odd}}
         =\ind_{k\submask n}.
  \label{eq:pascal-parity-matrix}
\end{equation}
As a locally finite lower-triangular matrix,
\begin{equation}
  \boxed{
  (P^{-1})_{n,k}=P_{n,k}\eps_{n\xor k}.}
  \label{eq:pascal-parity-inverse}
\end{equation}
In particular, the first column of \(P^{-1}\) is the signed Thue--Morse
sequence.  This is a second global Pascal encoding: \(\eps\) is not only a row
statistic but the M\"obius inverse of the entire parity incidence matrix.

Combining \eqref{eq:submask-reconstruction}, Lucas's indicator, and the central
binomial valuation gives a formula involving only ordinary binomial
coefficients:
\begin{theorem}[Nested binomial formula]
\label{thm:nested-binomial}
For every \(n\ge0\),
\begin{equation}
  \boxed{
  \eps_n=
  \sum_{k=0}^{n}
   \frac{1-(-1)^{\binom nk}}2
   (-2)^{\vtwo\binom{2k}{k}}.}
  \label{eq:nested-binomial}
\end{equation}
\end{theorem}

\begin{proof}
The first factor is \(\ind_{k\submask n}\) by
\eqref{eq:odd-indicator}, and the exponent in the second is \(\wt(k)\) by
\cref{thm:central-binomial}.  The sum is therefore exactly
\eqref{eq:submask-reconstruction}.
\end{proof}

\section[Generating products]{Finite and infinite generating products}
\label{sec:generating}

\subsection{The Boolean-cube mother identity}

A multivariate identity contains the finite generating product, the weight
enumerator, and the Walsh character formula as specializations.

\begin{theorem}[Boolean-cube product]
\label{thm:boolean-mother}
For commuting indeterminates \(x_0,\ldots,x_{m-1}\),
\begin{equation}
  \boxed{
  \sum_{n=0}^{2^m-1}\eps_n
     \prod_{j=0}^{m-1}x_j^{b_j(n)}
   =\prod_{j=0}^{m-1}(1-x_j).}
  \label{eq:boolean-mother}
\end{equation}
\end{theorem}

\begin{proof}
Identify \(n\) with its bit vector \(b\in\{0,1\}^m\).  The sum factors over
coordinates:
\[
 \sum_{b\in\{0,1\}^m}\prod_j(-x_j)^{b_j}
 =\prod_j\bigl(1-x_j\bigr).
\]
\end{proof}

Setting all \(x_j=-u\) recovers \eqref{eq:weight-enumerator}; setting
\(x_j=z^{2^j}\) gives the central finite polynomial.

\subsection{The finite Prouhet polynomial}

Define
\begin{equation}
  P_m(z)=\sum_{n=0}^{2^m-1}\eps_nz^n.
  \label{eq:Pm-definition}
\end{equation}

\begin{theorem}[Complete finite factorization]
\label{thm:finite-product}
For every \(m\ge0\),
\begin{equation}
  \boxed{
  P_m(z)=\prod_{j=0}^{m-1}(1-z^{2^j}).}
  \label{eq:finite-product}
\end{equation}
The empty product for \(m=0\) is \(1\).
\end{theorem}

\begin{proof}
Use \cref{thm:boolean-mother} with \(x_j=z^{2^j}\).  Equivalently, expanding
the product chooses a subset of powers of two; its sum is an integer
\(n<2^m\), uniquely, and the coefficient is
\((-1)^{\#\text{chosen bits}}=\eps_n\).
\end{proof}

\begin{corollary}[Coefficient and block formulae]
\label{cor:coefficient-block}
For every \(n\ge0\),
\begin{equation}
  \boxed{
  \eps_n=[z^n]\prod_{j\ge0}(1-z^{2^j}),}
  \label{eq:coefficient-formula}
\end{equation}
where only factors with \(2^j\le n\) affect the coefficient.  Moreover, for
\(q,m\ge0\),
\begin{equation}
  \sum_{r=0}^{2^m-1}\eps_{2^mq+r}z^r
  =\eps_qP_m(z).
  \label{eq:block-polynomial}
\end{equation}
\end{corollary}

\begin{proof}
The first statement is coefficient stabilization in the finite products.  The
second is \cref{thm:block-concatenation} summed over \(r\).
\end{proof}

\subsection{Reciprocity and cyclotomic multiplicities}

The reflection symmetry becomes a reciprocal-polynomial identity.

\begin{proposition}[Reciprocity]
\label{prop:reciprocity}
For every \(m\ge0\),
\begin{equation}
  z^{2^m-1}P_m(z^{-1})=(-1)^mP_m(z).
  \label{eq:Pm-reciprocity}
\end{equation}
Thus \(P_m\) is reciprocal for even \(m\) and anti-reciprocal for odd \(m\).
\end{proposition}

\begin{proof}
Reverse the coefficients and apply \eqref{eq:block-reflection}; alternatively,
apply \(z^{2^j}(1-z^{-2^j})=z^{2^j}-1=-(1-z^{2^j})\) to every factor.
\end{proof}

Repeatedly factor \(1-z^{2^j}\) by difference of squares:
\begin{equation}
  P_m(z)=(1-z)^m
   \prod_{q=0}^{m-2}(1+z^{2^q})^{m-1-q}.
  \label{eq:Pm-difference-squares}
\end{equation}
If \(\Phi_d\) is the \(d\)-th cyclotomic polynomial, then equivalently
\begin{equation}
  \boxed{
  P_m(z)=(-1)^m
   \prod_{q=0}^{m-1}\Phi_{2^q}(z)^{m-q}.}
  \label{eq:Pm-cyclotomic}
\end{equation}
Here \(\Phi_1(z)=z-1\).  In particular, a primitive \(2^q\)-th root of unity
is a zero of multiplicity \(m-q\) when \(0\le q<m\).  This multiplicity table
simultaneously explains Prouhet cancellation at \(z=1\) and the systematic
zeros in the dyadic discrete Fourier transform.

\subsection{The infinite product and Mahler equation}

For \(|z|<1\), define
\begin{equation}
  E(z)=\sum_{n\ge0}\eps_nz^n,
  \qquad
  T(z)=\sum_{n\ge0}\tau(n)z^n.
  \label{eq:E-T-definitions}
\end{equation}

\begin{theorem}[Infinite lacunary product]
\label{thm:infinite-product}
For \(|z|<1\), and coefficientwise as a formal power series,
\begin{equation}
  \boxed{
  E(z)=\prod_{j=0}^{\infty}(1-z^{2^j}).}
  \label{eq:infinite-product}
\end{equation}
The product converges uniformly on compact subsets of the unit disk.  The
zero--one series is
\begin{equation}
  \boxed{
  T(z)=\frac{1}{2(1-z)}-\frac12E(z).}
  \label{eq:T-from-E}
\end{equation}
\end{theorem}

\begin{proof}
For \(|z|\le\rho<1\), the sum \(\sum_j|z|^{2^j}\) converges, so the product
converges uniformly on the compact disk.  Every coefficient stabilizes in the
finite products \(P_m\), where \cref{thm:finite-product} identifies it as
\(\eps_n\).  The formula for \(T\) follows from
\(\tau(n)=(1-\eps_n)/2\).
\end{proof}

Peeling off factors gives the Mahler and dissection equations.

\begin{corollary}[Mahler equations and dissections]
\label{cor:mahler-equations}
For \(|z|<1\),
\begin{align}
  E(z)&=(1-z)E(z^2),
  \label{eq:E-Mahler}\\
  E(z)&=P_m(z)E(z^{2^m}),
  \label{eq:E-dissection}\\
  T(z)&=(1-z)T(z^2)+\frac{z}{1-z^2}.
  \label{eq:T-Mahler}
\end{align}
\end{corollary}

\begin{proof}
Separate the first factor, or the first \(m\) factors, in
\eqref{eq:infinite-product}.  Substitute \eqref{eq:T-from-E} for the final
identity; it also follows directly from \eqref{eq:tau-recursion}.
\end{proof}

Iterating the inhomogeneous bit equation yields a useful lacunary sum.

\begin{theorem}[Lacunary quotient expansion]
\label{thm:lacunary-quotient}
Let \(P_0(z)=1\) and \(P_j(z)=\prod_{h=0}^{j-1}(1-z^{2^h})\) for \(j\ge1\).
Then, for \(|z|<1\),
\begin{equation}
  \boxed{
  T(z)=\sum_{j=0}^{\infty}
    P_j(z)\frac{z^{2^j}}{1-z^{2^{j+1}}}.}
  \label{eq:lacunary-quotient}
\end{equation}
\end{theorem}

\begin{proof}
After \(M\) iterations of \eqref{eq:T-Mahler},
\[
 T(z)=\sum_{j=0}^{M-1}P_j(z)
       \frac{z^{2^j}}{1-z^{2^{j+1}}}
       +P_M(z)T(z^{2^M}).
\]
On compact subdisks \(P_M\) is bounded and \(T(z^{2^M})\to T(0)=0\).
\end{proof}

\subsection{Thue--Morse constants in every base}

For a real \(b>1\), put
\begin{equation}
  \Theta_b=\sum_{n\ge0}\frac{\tau(n)}{b^{n+1}}.
  \label{eq:base-b-constant-definition}
\end{equation}
For integer \(b\ge2\), these are the base-\(b\) numbers whose fractional
digits are the Thue--Morse bits.

\begin{corollary}[Base-\(b\) product]
\label{cor:base-b-product}
For every real \(b>1\),
\begin{equation}
  \boxed{
  \Theta_b=\frac{1}{2(b-1)}-
   \frac1{2b}\prod_{j\ge0}\left(1-b^{-2^j}\right).}
  \label{eq:base-b-product}
\end{equation}
The signed companion is
\begin{equation}
  \sum_{n\ge0}\frac{\eps_n}{b^{n+1}}
   =\frac1b\prod_{j\ge0}\left(1-b^{-2^j}\right).
  \label{eq:signed-base-b}
\end{equation}
\end{corollary}

\begin{proof}
Set \(z=1/b\) in \eqref{eq:T-from-E} and
\eqref{eq:infinite-product}, then multiply by \(1/b\).
\end{proof}

For \(b=2\), \(\Theta_2=0.0110100110010110\ldots{}_2\) is the classical
Prouhet--Thue--Morse constant.  Mahler's theory implies its transcendence; that
deep arithmetic result is separate from the elementary product identity
proved here \cite{mahler1929,alloucheShallit2003}.

\section[Euler transforms and algebraic series]{Euler transforms, coefficient recurrences, and algebraic series}
\label{sec:euler}

\subsection{The logarithmic ruler transform}

Take the formal logarithm of \eqref{eq:infinite-product}:
\begin{align}
 \log E(z)
 &=\sum_{j\ge0}\log(1-z^{2^j})
  =-\sum_{j\ge0}\sum_{r\ge1}\frac{z^{r2^j}}r.
  \label{eq:log-E-double}
\end{align}
For a fixed exponent \(k\), the admissible \(j\)'s satisfy
\(0\le j\le\vtwo(k)\), and the corresponding denominator is \(k/2^j\).
Define
\begin{equation}
  a_k=\sum_{j=0}^{\vtwo(k)}2^j
      =2^{\vtwo(k)+1}-1.
  \label{eq:ak-definition}
\end{equation}

\begin{theorem}[Euler transform of the ruler function]
\label{thm:log-E}
As a formal series, and analytically for \(|z|<1\),
\begin{equation}
  \boxed{
  \log E(z)=-\sum_{k\ge1}\frac{a_k}{k}z^k,
  \qquad a_k=2^{\vtwo(k)+1}-1.}
  \label{eq:log-E}
\end{equation}
\end{theorem}

This representation contains no explicit binary digit sum.  The coefficient
sequence \((a_k)\) is a simple transform of the ruler sequence.

\subsection{A nonlocal convolution recurrence}

\begin{theorem}[Ruler-convolution recurrence]
\label{thm:ruler-convolution}
With \(\eps_0=1\), for every \(n\ge1\),
\begin{equation}
  \boxed{
  n\eps_n=-\sum_{k=1}^{n}a_k\eps_{n-k}.}
  \label{eq:ruler-convolution}
\end{equation}
\end{theorem}

\begin{proof}
Differentiate \eqref{eq:log-E}, multiply by \(E(z)\), and compare the
coefficient of \(z^{n-1}\):
\[
 E'(z)=-E(z)\sum_{k\ge1}a_kz^{k-1}.
\]
\end{proof}

The weighted convolution on the right always collapses to \(\pm n\).  As an
evaluator it is inferior to popcount, but it is the natural recurrence attached
to the Euler product.

\subsection{Partition and Bell-polynomial formulae}

Exponentiating \eqref{eq:log-E} gives a finite partition sum for every
coefficient.

\begin{theorem}[Partition formula]
\label{thm:partition-formula}
For \(n\ge0\),
\begin{equation}
  \boxed{
  \eps_n=
  \sum_{\substack{m_1,\ldots,m_n\ge0\\
        m_1+2m_2+\cdots+nm_n=n}}
  \prod_{j=1}^{n}
   \frac1{m_j!}\left(-\frac{a_j}{j}\right)^{m_j}.}
  \label{eq:partition-formula}
\end{equation}
\end{theorem}

\begin{proof}
Write
\(E(z)=\prod_{j\ge1}\exp(-a_jz^j/j)\), expand each exponential, and collect
terms of total degree \(n\).
\end{proof}

Let \(\Bell_n\) denote the complete exponential Bell polynomial, defined by
\begin{equation}
  \exp\!\left(\sum_{j\ge1}x_j\frac{z^j}{j!}\right)
  =\sum_{n\ge0}\Bell_n(x_1,\ldots,x_n)\frac{z^n}{n!}.
  \label{eq:Bell-definition}
\end{equation}

\begin{corollary}[Bell-polynomial coefficient]
\label{cor:Bell-coefficient}
For every \(n\ge0\),
\begin{equation}
  \boxed{
  \eps_n=\frac1{n!}
   \Bell_n\bigl(-0!a_1,-1!a_2,\ldots,-(n-1)!a_n\bigr).}
  \label{eq:Bell-coefficient}
\end{equation}
\end{corollary}

\begin{proof}
Use \(x_j=-(j-1)!a_j\) in \eqref{eq:Bell-definition}.
\end{proof}

Although the right sides are rational-looking sums, they always equal
\(\pm1\).  This integrality collapse is itself encoded by the original
lacunary product.

\subsection{A Hessenberg determinant}

For \(n\ge1\), let
\begin{equation}
 H_n=
 \begin{pmatrix}
 a_1&1&0&\cdots&0\\
 a_2&a_1&2&\ddots&\vdots\\
 a_3&a_2&a_1&\ddots&0\\
 \vdots&\vdots&\ddots&\ddots&n-1\\
 a_n&a_{n-1}&\cdots&a_2&a_1
 \end{pmatrix}.
 \label{eq:Hessenberg-matrix}
\end{equation}

\begin{theorem}[Determinant formula]
\label{thm:Hessenberg-determinant}
For every \(n\ge1\),
\begin{equation}
  \boxed{
  \eps_n=\frac{(-1)^n}{n!}\det H_n.}
  \label{eq:Hessenberg-determinant}
\end{equation}
In particular,
\begin{equation}
  \det H_n=(-1)^n n!\eps_n\in\{n!,-n!\}.
  \label{eq:Hessenberg-factorial}
\end{equation}
\end{theorem}

\begin{proof}
Put \(d_0=1\) and \(d_n=\det H_n\).  Expansion of a lower Hessenberg
determinant along its first column gives
\begin{equation}
 d_n=\sum_{k=1}^{n}(-1)^{k-1}
     \frac{(n-1)!}{(n-k)!}a_kd_{n-k}.
 \label{eq:Hessenberg-recurrence}
\end{equation}
Set \(d_n=(-1)^nn!c_n\).  Then \eqref{eq:Hessenberg-recurrence} reduces to
\[
 nc_n=-\sum_{k=1}^{n}a_kc_{n-k},
 \]
with \(c_0=1\).  By uniqueness of the recurrence in
\cref{thm:ruler-convolution}, \(c_n=\eps_n\).
\end{proof}

\subsection{The zero--one series over \texorpdfstring{\(\mathbb F_2\)}{F2}}
\label{sec:F2-series}

The signed series loses all information modulo two because \(-1=1\).  The bit
series remains nontrivial.  Work in \(\Ftwo[[z]]\) and write
\begin{equation}
  \Theta(z)=\sum_{n\ge0}\tau(n)z^n.
  \label{eq:Theta-F2}
\end{equation}

\begin{theorem}[Artin--Schreier equation]
\label{thm:Artin-Schreier}
In \(\Ftwo[[z]]\),
\begin{align}
  \Theta(z)&=(1+z)\Theta(z^2)+\frac{z}{1+z^2},
  \label{eq:Theta-functional}\\
  \boxed{
  (1+z)^3\Theta(z)^2+(1+z)^2\Theta(z)+z=0.}
  \label{eq:Theta-quadratic}
\end{align}
If \(Y(z)=(1+z)\Theta(z)\), then
\begin{equation}
  \boxed{
  Y(z)^2+Y(z)=\frac{z}{1+z}.}
  \label{eq:Y-Artin-Schreier}
\end{equation}
\end{theorem}

\begin{proof}
Split the series into even and odd coefficients and use
\(\tau(2n)=\tau(n)\), \(\tau(2n+1)=1+\tau(n)\) in characteristic two:
\[
 \Theta(z)=\Theta(z^2)+z\sum_{n\ge0}(1+\tau(n))z^{2n}
 =(1+z)\Theta(z^2)+\frac{z}{1+z^2}.
\]
Frobenius gives \(\Theta(z^2)=\Theta(z)^2\), while
\(1+z^2=(1+z)^2\).  Clear the denominator and substitute
\(Y=(1+z)\Theta\).
\end{proof}

\begin{corollary}[Explicit Artin--Schreier sum]
\label{cor:Artin-Schreier-sum}
The unique zero-constant-term solution is
\begin{align}
  Y(z)&=\sum_{j=0}^{\infty}
      \left(\frac{z}{1+z}\right)^{2^j}
      =\sum_{j=0}^{\infty}\frac{z^{2^j}}{1+z^{2^j}},
  \label{eq:Y-explicit}\\
  \boxed{
  \Theta(z)=\frac1{1+z}
      \sum_{j=0}^{\infty}\frac{z^{2^j}}{1+z^{2^j}}.}
  \label{eq:Theta-explicit}
\end{align}
\end{corollary}

\begin{proof}
Let \(A=z/(1+z)\).  In the \(z\)-adic topology,
\(Y=\sum_{j\ge0}A^{2^j}\) is well defined and
\(Y^2+Y=A\) by index shifting.  Frobenius gives
\(A^{2^j}=z^{2^j}/(1+z^{2^j})\).  A zero-constant solution of the homogeneous
equation \(Z^2+Z=0\) must vanish, proving uniqueness.
\end{proof}

Christol's theorem says that a finite-field sequence is automatic exactly when
its generating series is algebraic.  Here the degree-two certificate is
derived directly; it is an algebraic shadow of the two-state automaton
\cite{christol1980}.

\subsection{An integral algebraic lift}
\label{sec:integer-lift}

The fourth draft records a rather different integer sequence whose parity is
Thue--Morse.  Define \(c_0=0\), \(c_1=1\), and for \(n\ge2\),
\begin{equation}
  nc_n=(5-2n)c_{n-1}+3(n-1)c_{n-2}+1.
  \label{eq:integer-lift-recurrence}
\end{equation}
Its first terms are
\[
 0,1,1,2,1,4,-2,13,-23,68,-164,439,-1146,3067,\ldots .
\]

\begin{theorem}[Integer algebraic lift]
\label{thm:integer-lift}
Let \(C(z)=\sum_{n\ge0}c_nz^n\).  Then
\begin{align}
  C(z)&=\frac{1}{2(1-z)}
   \left(\sqrt{\frac{1+3z}{1-z}}-1\right),
  \label{eq:integer-lift-closed}\\
  (1-z)^3C(z)^2+(1-z)^2C(z)&=z.
  \label{eq:integer-lift-algebraic}
\end{align}
All \(c_n\) are integers and
\begin{equation}
  \boxed{
  c_n\equiv\tau(n)\pmod2,
  \qquad
  \eps_n=(-1)^{c_n}.}
  \label{eq:integer-lift-parity}
\end{equation}
\end{theorem}

\begin{proof}
Multiplying \eqref{eq:integer-lift-recurrence} by \(z^{n-1}\), summing, and
including the initial terms gives a first-order differential equation whose
zero-constant solution is \eqref{eq:integer-lift-closed}; squaring yields
\eqref{eq:integer-lift-algebraic}.

Conversely, the coefficient of \(z^n\) in \((1-z)^2C(z)\) contains \(c_n\)
with coefficient one, while the quadratic term uses only earlier coefficients
because \(C(0)=0\).  Thus \eqref{eq:integer-lift-algebraic} recursively
determines integers \(c_n\).

Modulo two, \eqref{eq:integer-lift-algebraic} becomes precisely
\eqref{eq:Theta-quadratic}.  That equation has a unique zero-constant solution
in \(\Ftwo[[z]]\), so \(C(z)\equiv\Theta(z)\pmod2\).
\end{proof}

\begin{remark}
The lift grows and changes sign; it is not a disguised digit sum.  Its parity
is Thue--Morse because an integral algebraic equation reduces to the automatic
sequence's Artin--Schreier equation.  This suggests a reusable discovery
method: search for integral algebraic lifts of finite-field automatic series.
\end{remark}

\section[Finite differences and moments]{Finite differences, Prouhet cancellation, and all dyadic moments}
\label{sec:finite-differences}

\subsection{The mixed-difference master identity}

For a function \(f\), let
\begin{equation}
  (\Delta_hf)(x)=f(x+h)-f(x),
  \qquad
  (T_hf)(x)=f(x+h).
  \label{eq:difference-shift}
\end{equation}
Thus \(I-T_h=-\Delta_h\).

\begin{theorem}[Dyadic finite-difference factorization]
\label{thm:finite-difference-master}
For every \(m\ge0\), every scalar \(h\), and every function for which the
values are defined,
\begin{align}
  \boxed{
  \sum_{n=0}^{2^m-1}\eps_nf(x+nh)
  =\prod_{j=0}^{m-1}(I-T_{2^jh})f(x)}
  \label{eq:operator-product}\\
  \boxed{
  \phantom{\sum_{n=0}^{2^m-1}\eps_nf(x+nh)}
  =(-1)^m
    (\Delta_h\Delta_{2h}\cdots\Delta_{2^{m-1}h}f)(x).}
  \label{eq:mixed-difference}
\end{align}
\end{theorem}

\begin{proof}
Expand the operator product.  A subset \(S\subseteq\{0,\ldots,m-1\}\)
contributes the shift by \(h\sum_{j\in S}2^j\) and sign \((-1)^{|S|}\).
Binary expansion identifies the subset sum with a unique \(n<2^m\), and the
sign is \(\eps_n\).
\end{proof}

This theorem is stronger than any one moment identity: it is an
inclusion--exclusion formula valid for arbitrary test functions.

\subsection{Exact multiple-integral remainder}

\begin{corollary}[Integral cubature formula]
\label{cor:integral-cubature}
Suppose \(h\ge0\) and \(f\in C^m\) on an interval containing all relevant
points.  Then
\begin{align}
 \sum_{n<2^m}\eps_nf(x+nh)
 &=(-1)^m
 \int_0^h\int_0^{2h}\cdots\int_0^{2^{m-1}h}
 f^{(m)}(x+t_0+\cdots+t_{m-1})
 \dd t_{m-1}\cdots\dd t_0
 \label{eq:integral-cubature}\\
 &=(-1)^m2^{\binom m2}h^m
 \int_{[0,1]^m}
 f^{(m)}\!\left(x+h\sum_{j=0}^{m-1}2^ju_j\right)
 \dd\bm u.
 \label{eq:unit-cube-cubature}
\end{align}
Consequently
\begin{equation}
  \left|\sum_{n<2^m}\eps_nf(x+nh)\right|
  \le2^{\binom m2}h^m
  \sup_{x\le y\le x+(2^m-1)h}|f^{(m)}(y)|.
  \label{eq:finite-difference-bound}
\end{equation}
\end{corollary}

\begin{proof}
Apply
\((I-T_a)f(x)=-\int_0^af'(x+t)\dd t\) successively.  Rescale
\(t_j=2^jhu_j\) for the unit-cube form.  The box volume is
\(h^m2^{0+1+\cdots+(m-1)}\).
\end{proof}

\begin{corollary}[Sign for completely monotone kernels]
\label{cor:complete-monotone}
If \((-1)^mf^{(m)}\ge0\) on the interval, then
\begin{equation}
  \sum_{n<2^m}\eps_nf(x+nh)\ge0.
  \label{eq:complete-monotone-sign}
\end{equation}
For \(a,s>0\),
\begin{equation}
  \boxed{
  \sum_{n<2^m}\frac{\eps_n}{(a+n)^s}
  =(s)_m
  \int_0^1\int_0^2\cdots\int_0^{2^{m-1}}
  \frac{\dd t_{m-1}\cdots\dd t_0}
       {(a+t_0+\cdots+t_{m-1})^{s+m}}>0,}
  \label{eq:reciprocal-power-positive}
\end{equation}
where \((s)_m=s(s+1)\cdots(s+m-1)\).
For \(m\ge1\),
\begin{equation}
  \sum_{n<2^m}\eps_n\log(a+n)<0,
  \qquad
  0<\prod_{n<2^m}(a+n)^{\eps_n}<1.
  \label{eq:log-product-sign}
\end{equation}
\end{corollary}

\begin{proof}
Use \(f(y)=(a+y)^{-s}\), for which
\(f^{(m)}(y)=(-1)^m(s)_m(a+y)^{-s-m}\).  For the logarithm,
\(f^{(m)}(y)=(-1)^{m-1}(m-1)!/(a+y)^m\), so the integral is strictly negative.
Exponentiation gives the product inequality.
\end{proof}

\subsection{Prouhet cancellation and the first surviving moment}

\begin{theorem}[Prouhet polynomial annihilation]
\label{thm:prouhet}
If \(p\) is a polynomial of degree less than \(m\), then
\begin{equation}
  \boxed{
  \sum_{n=0}^{2^m-1}\eps_np(x+nh)=0.}
  \label{eq:prouhet-affine}
\end{equation}
If \(p(y)=cy^m+\) lower-degree terms, then
\begin{equation}
  \boxed{
  \sum_{n=0}^{2^m-1}\eps_np(x+nh)
  =(-1)^mm!2^{\binom m2}h^m c.}
  \label{eq:prouhet-sharp-polynomial}
\end{equation}
\end{theorem}

\begin{proof}
Every \(m\)-fold finite difference annihilates a polynomial of degree below
\(m\).  At degree exactly \(m\), its \(m\)-th derivative is the constant
\(m!c\), so \eqref{eq:unit-cube-cubature} gives the sharp value.
\end{proof}

For \(h=1\), let
\begin{equation}
  A_m=\{0\le n<2^m:\tau(n)=0\},
  \qquad
  B_m=\{0\le n<2^m:\tau(n)=1\}.
  \label{eq:prouhet-sets}
\end{equation}
Then
\begin{equation}
  \sum_{n\in A_m}n^r=\sum_{n\in B_m}n^r,
  \qquad 0\le r<m,
  \label{eq:prouhet-partition}
\end{equation}
which is Prouhet's classical equal-sums-of-like-powers construction.  At
\(r=m\), the difference is nonzero, so the cancellation order is exact.

\subsection{The finite exponential transform}

Define
\begin{equation}
  M_{m,r}=\sum_{n=0}^{2^m-1}\eps_nn^r.
  \label{eq:moment-definition}
\end{equation}
Substitute \(z=\e^t\) in \eqref{eq:finite-product}:
\begin{equation}
  \boxed{
  \sum_{r\ge0}M_{m,r}\frac{t^r}{r!}
  =\sum_{n<2^m}\eps_n\e^{nt}
  =\prod_{j=0}^{m-1}(1-\e^{2^jt}).}
  \label{eq:moment-egf}
\end{equation}
The order-\(m\) zero at \(t=0\) immediately gives
\begin{equation}
  M_{m,r}=0\quad(r<m),
  \qquad
  \boxed{M_{m,m}=(-1)^mm!2^{\binom m2}.}
  \label{eq:first-surviving-moment}
\end{equation}

\subsection{All power moments: a finite composition formula}

\begin{theorem}[Complete moment composition]
\label{thm:moment-composition}
For \(m\ge1\) and \(r\ge m\),
\begin{equation}
  \boxed{
  M_{m,r}=(-1)^m r!\,2^{\binom m2}
  \sum_{\substack{q_0,\ldots,q_{m-1}\ge0\\
                  q_0+\cdots+q_{m-1}=r-m}}
  \prod_{j=0}^{m-1}
    \frac{2^{jq_j}}{(q_j+1)!}.}
  \label{eq:moment-composition}
\end{equation}
Equivalently,
\begin{equation}
  M_{m,r}=(-1)^m r!
  \sum_{\substack{p_0,\ldots,p_{m-1}\ge1\\
                  p_0+\cdots+p_{m-1}=r}}
  \frac{2^{\sum_{j=0}^{m-1}jp_j}}
       {p_0!\cdots p_{m-1}!}.
  \label{eq:moment-positive-compositions}
\end{equation}
\end{theorem}

\begin{proof}
Factor the lowest-order term from each factor in \eqref{eq:moment-egf}:
\begin{align*}
 \prod_{j=0}^{m-1}(1-\e^{2^jt})
 &=(-1)^m2^{\binom m2}t^m
   \prod_{j=0}^{m-1}
    \frac{\e^{2^jt}-1}{2^jt}\\
 &=(-1)^m2^{\binom m2}t^m
   \prod_{j=0}^{m-1}
    \sum_{q\ge0}\frac{2^{jq}t^q}{(q+1)!}.
\end{align*}
Extract \(t^r\) and multiply by \(r!\).  Set \(p_j=q_j+1\) for the second
form.
\end{proof}

\subsection{Bell-polynomial compression and cumulants}

Let \(B_{2\ell}\) be the Bernoulli numbers with \(B_2=1/6\), and define
\begin{align}
 \kappa_1^{(m)}&=\frac{2^m-1}{2},
 \label{eq:kappa-one}\\
 \kappa_{2\ell}^{(m)}
 &=\frac{B_{2\ell}}{2\ell}
   \frac{2^{2\ell m}-1}{2^{2\ell}-1}
   \qquad(\ell\ge1),
 \label{eq:kappa-even}\\
 \kappa_{2\ell+1}^{(m)}&=0
   \qquad(\ell\ge1).
 \label{eq:kappa-odd}
\end{align}

\begin{theorem}[Bell-polynomial moments]
\label{thm:Bell-moments}
For \(m\ge1\) and \(r\ge0\),
\begin{equation}
  \boxed{
  M_{m,m+r}=(-1)^m2^{\binom m2}
    \frac{(m+r)!}{r!}
    \Bell_r\!\left(
       \kappa_1^{(m)},\ldots,\kappa_r^{(m)}
    \right).}
  \label{eq:Bell-moments}
\end{equation}
\end{theorem}

\begin{proof}
The classical formal expansion
\begin{equation}
  \log\!\left(\frac{\e^y-1}{y}\right)
  =\frac y2+
   \sum_{\ell\ge1}
   \frac{B_{2\ell}}{2\ell(2\ell)!}y^{2\ell}
  \label{eq:log-expm1}
\end{equation}
turns the normalized product into
\begin{equation}
 \prod_{j=0}^{m-1}(1-\e^{2^jt})
 =(-1)^m2^{\binom m2}t^m
  \exp\!\left(
    \sum_{r\ge1}\kappa_r^{(m)}\frac{t^r}{r!}
  \right).
  \label{eq:moment-Bell-factor}
\end{equation}
Use the defining generating function \eqref{eq:Bell-definition} and compare
coefficients.
\end{proof}

The first three nonzero layers are
\begin{align}
 M_{m,m}
 &=(-1)^m2^{\binom m2}m!,
 \label{eq:moment-layer0}\\
 M_{m,m+1}
 &=(-1)^m2^{\binom m2}(m+1)!\frac{2^m-1}{2},
 \label{eq:moment-layer1}\\
 M_{m,m+2}
 &=(-1)^m2^{\binom m2}(m+2)!
   \left(
     \frac{(2^m-1)^2}{8}+\frac{4^m-1}{72}
   \right).
 \label{eq:moment-layer2}
\end{align}

\subsection{Centered moments and a parity selection rule}
\label{sec:centered-moments}

Let
\begin{equation}
  c_m=\frac{2^m-1}{2}.
  \label{eq:block-center}
\end{equation}
The centered exponential transform is especially symmetric.

\begin{theorem}[Centered sinh product]
\label{thm:centered-sinh}
For every \(m\ge0\),
\begin{equation}
  \boxed{
  \sum_{n=0}^{2^m-1}\eps_n\e^{(n-c_m)t}
  =(-2)^m\prod_{j=0}^{m-1}\sinh(2^{j-1}t).}
  \label{eq:centered-sinh}
\end{equation}
Consequently
\begin{equation}
  \boxed{
  \sum_{n<2^m}\eps_n(n-c_m)^r=0}
  \label{eq:centered-parity-zero}
\end{equation}
whenever \(r<m\) or \(r\not\equiv m\pmod2\).  At \(r=m\),
\begin{equation}
  \sum_{n<2^m}\eps_n(n-c_m)^m
  =(-1)^mm!2^{\binom m2}.
  \label{eq:centered-first}
\end{equation}
\end{theorem}

\begin{proof}
Use
\(1-\e^u=-2\e^{u/2}\sinh(u/2)\) in
\eqref{eq:moment-egf}.  The product of exponential factors is
\(\e^{c_mt}\), which disappears after centering.  The right side is a product
of \(m\) odd functions, hence has parity \(m\), and each factor starts with a
linear term.  Coefficient comparison proves the selection rule and the first
coefficient.
\end{proof}

This extra parity cancellation is invisible in the uncentered Prouhet
statement.  It is a consequence of combining dyadic reflection with moment
cancellation.

\subsection{All binomial-basis moments}

Define
\begin{equation}
  B_{m,r}=\sum_{n=0}^{2^m-1}\eps_n\binom nr.
  \label{eq:binomial-moment-definition}
\end{equation}

\begin{theorem}[Binomial moment product]
\label{thm:binomial-moments}
For \(\ell\ge0\),
\begin{equation}
  \boxed{
  B_{m,m+\ell}=(-1)^m2^{\binom m2}
  [u^\ell]\prod_{j=0}^{m-1}
  \frac{(1+u)^{2^j}-1}{2^ju}.}
  \label{eq:binomial-moment-product}
\end{equation}
For \(r<m\), \(B_{m,r}=0\); in particular,
\begin{equation}
  B_{m,m}=(-1)^m2^{\binom m2}.
  \label{eq:first-binomial-moment}
\end{equation}
\end{theorem}

\begin{proof}
Expand \(P_m(1+u)\) in two ways:
\[
 P_m(1+u)=\sum_{n<2^m}\eps_n(1+u)^n
 =\sum_{r\ge0}B_{m,r}u^r,
\]
while
\[
 P_m(1+u)=(-1)^m2^{\binom m2}u^m
 \prod_{j=0}^{m-1}\frac{(1+u)^{2^j}-1}{2^ju}.
\]
Compare coefficients.
\end{proof}

The binomial basis is often preferable in discrete formalization because
finite differences obey the simple lowering rule
\(\Delta\binom xr=\binom{x}{r-1}\).

\section[Sparse Prouhet identities]{Sparse Prouhet identities on arbitrary Pascal supports}
\label{sec:sparse-prouhet}

The ordinary Prouhet block \([0,2^m)\) is the submask lattice of
\(2^m-1\).  Lucas's theorem suggests replacing this all-ones mask by an
arbitrary \(n\).

Define
\begin{equation}
  J(n)=\{j\ge0:b_j(n)=1\},
  \qquad
  s=|J(n)|=\wt(n),
  \qquad
  \beta(n)=\sum_{j\in J(n)}j.
  \label{eq:J-beta}
\end{equation}

\begin{theorem}[Signed Pascal-support product]
\label{thm:signed-Pascal-support}
For every \(n\ge0\),
\begin{equation}
  \boxed{
  R_n(z):=\sum_{k\submask n}\eps_kz^k
  =\prod_{j\in J(n)}(1-z^{2^j}).}
  \label{eq:signed-Pascal-support}
\end{equation}
Equivalently,
\begin{equation}
  R_n(z)=\sum_{k=0}^{n}
  \eps_k\ind_{\binom nk\text{ odd}}z^k.
  \label{eq:signed-Pascal-Lucas}
\end{equation}
\end{theorem}

\begin{proof}
Expanding the product chooses a subset of the occupied bit positions of
\(n\).  Its exponent is the corresponding submask \(k\), and its sign is
\((-1)^{\wt(k)}=\eps_k\).  Lucas's criterion gives the second form.
\end{proof}

\begin{corollary}[Balance on odd Pascal positions]
\label{cor:Pascal-support-balance}
For every \(n>0\),
\begin{equation}
  \sum_{\substack{0\le k\le n\\\binom nk\text{ odd}}}\eps_k=0.
  \label{eq:Pascal-support-balance}
\end{equation}
Thus exactly half of the odd positions in every nonzero Pascal row have even
binary weight and half have odd binary weight.
\end{corollary}

\begin{proof}
Set \(z=1\) in \eqref{eq:signed-Pascal-support}.  There are
\(2^{\wt(n)}\) support positions by \cref{cor:glaisher-count}.
\end{proof}

\begin{theorem}[Sparse Prouhet theorem]
\label{thm:sparse-Prouhet}
Let \(n>0\), \(s=\wt(n)\).  For every polynomial \(p\) of degree below
\(s\),
\begin{equation}
  \boxed{
  \sum_{k\submask n}\eps_kp(k)=0.}
  \label{eq:sparse-Prouhet}
\end{equation}
The first surviving power moment is
\begin{equation}
  \boxed{
  \sum_{k\submask n}\eps_kk^s
   =(-1)^ss!2^{\beta(n)}.}
  \label{eq:sparse-first-moment}
\end{equation}
\end{theorem}

\begin{proof}
Every factor in \eqref{eq:signed-Pascal-support} has a simple zero at
\(z=1\), so \(R_n\) has a zero of exact order \(s\), with local leading term
\[
 R_n(z)=(-1)^s2^{\beta(n)}(z-1)^s+\bigO((z-1)^{s+1}).
\]
The derivative/falling-factorial proof of \cref{thm:prouhet} applies with the
sum restricted to submasks.
\end{proof}

\begin{corollary}[Sparse operator identity]
\label{cor:sparse-operator}
For every function \(f\),
\begin{equation}
  \boxed{
  \sum_{k\submask n}\eps_kf(x+hk)
   =\prod_{j\in J(n)}(I-T_{2^jh})f(x).}
  \label{eq:sparse-operator}
\end{equation}
If \(f\in C^s\) and \(h\ge0\), its absolute value is at most
\begin{equation}
  2^{\beta(n)}h^s
  \sup_{x\le y\le x+nh}|f^{(s)}(y)|.
  \label{eq:sparse-operator-bound}
\end{equation}
\end{corollary}

\begin{proof}
Expand the product over the set-bit positions and apply the fundamental
theorem of calculus once for each step.
\end{proof}

The whole sparse moment sequence also has a direct finite formula.

\begin{corollary}[All sparse moments]
\label{cor:all-sparse-moments}
For \(r\ge s\),
\begin{equation}
  \boxed{
  \sum_{k\submask n}\eps_kk^r
  =(-1)^sr!2^{\beta(n)}
  \sum_{\substack{(q_j)_{j\in J(n)}\ge0\\
                  \sum_{j\in J(n)}q_j=r-s}}
  \prod_{j\in J(n)}\frac{2^{jq_j}}{(q_j+1)!}.}
  \label{eq:all-sparse-moments}
\end{equation}
\end{corollary}

\begin{proof}
Substitute \(z=\e^t\) in \eqref{eq:signed-Pascal-support}, factor
\(-2^jt\) from every occupied bit, and extract coefficients exactly as in
\cref{thm:moment-composition}.
\end{proof}

\begin{example}
For \(n=13=(1101)_2\),
\(J(13)=\{0,2,3\}\), \(s=3\), and \(\beta(13)=5\).  The odd positions in
row \(13\) are
\[
 0,1,4,5,8,9,12,13.
\]
Their signed sums against \(1,k,k^2\) vanish, while
\[
 \sum_{k\submask13}\eps_kk^3=(-1)^3\,3!\,2^5=-192.
\]
This is a Prouhet partition on a sparse, Pascal-determined grid rather than a
consecutive interval.
\end{example}

\section[Prefix sums and the Fabius bridge]{Exact prefix sums, iterated antiderivatives, and the Fabius bridge}
\label{sec:prefix-fabius}

\subsection{The first prefix discrepancy}

Define the exclusive signed prefix and the number of one-bits in the sequence
prefix by
\begin{equation}
  S(N)=\sum_{n=0}^{N-1}\eps_n,
  \qquad
  A(N)=\sum_{n=0}^{N-1}\tau(n).
  \label{eq:prefix-definitions}
\end{equation}

\begin{theorem}[Exact prefix formula]
\label{thm:exact-prefix}
For every \(q\ge0\),
\begin{equation}
  \boxed{
  S(2q)=0,
  \qquad
  S(2q+1)=\eps_q.}
  \label{eq:exact-prefix}
\end{equation}
Equivalently,
\begin{equation}
  S(N)=\frac{1-(-1)^N}{2}\,
       \eps_{\lfloor N/2\rfloor}.
  \label{eq:compact-prefix}
\end{equation}
Therefore
\begin{equation}
  \boxed{
  A(N)=\frac{N-S(N)}2}
  \label{eq:ones-prefix}
\end{equation}
and explicitly
\begin{equation}
  A(2q)=q,
  \qquad
  A(2q+1)=q+\tau(q).
  \label{eq:ones-prefix-cases}
\end{equation}
\end{theorem}

\begin{proof}
Every adjacent pair cancels:
\(\eps_{2j}+\eps_{2j+1}=\eps_j-\eps_j=0\).  An odd-length prefix leaves the
unpaired term \(\eps_{2q}=\eps_q\).  The count formula follows from
\(\eps_n=1-2\tau(n)\).
\end{proof}

The inclusive sum through \(N\) is
\begin{equation}
  \sum_{n=0}^{N}\eps_n
  =\frac{1+(-1)^N}{2}\eps_{\lfloor N/2\rfloor}.
  \label{eq:inclusive-prefix}
\end{equation}

\begin{corollary}[Interval discrepancy]
\label{cor:interval-discrepancy}
For \(0\le a\le b\),
\begin{align}
  \sum_{n=a}^{b-1}\eps_n&=S(b)-S(a),
  \label{eq:interval-exact}\\
  \left|\sum_{n=a}^{b-1}\eps_n\right|&\le2,
  \label{eq:interval-signed-bound}\\
  \left|\sum_{n=a}^{b-1}\tau(n)-\frac{b-a}{2}\right|&\le1.
  \label{eq:interval-bit-bound}
\end{align}
Both constants are attained.
\end{corollary}

\begin{proof}
The interval sum is a difference of two values in \(\{-1,0,1\}\).  Convert
to zero--one values by \eqref{eq:normalization-bridge}.  For sharpness,
\(\eps_5+\eps_6=2\).
\end{proof}

\begin{corollary}[Aligned and partial dyadic blocks]
\label{cor:dyadic-block-prefix}
For \(m\ge1\) and \(q\ge0\),
\begin{equation}
  \sum_{r=0}^{2^m-1}\eps_{2^mq+r}=0.
  \label{eq:aligned-block-zero}
\end{equation}
More generally, for \(0\le R\le2^m\),
\begin{equation}
  \sum_{r=0}^{R-1}\eps_{2^mq+r}=\eps_qS(R).
  \label{eq:partial-block-prefix}
\end{equation}
\end{corollary}

\begin{proof}
Use \eqref{eq:block-sign} and factor out \(\eps_q\); then apply
\cref{thm:exact-prefix}.
\end{proof}

The prefix generating function is unusually simple:
\begin{equation}
  \sum_{N\ge0}S(N)z^N
  =\frac{z}{1-z}E(z)
  =zE(z^2)
  =z\prod_{j\ge1}(1-z^{2^j}).
  \label{eq:prefix-generating-first}
\end{equation}
The middle equality is the Mahler equation.

\subsection{Iterated inclusive prefixes}

Define discrete antiderivatives by
\begin{equation}
  S_n^{(0)}=\eps_n,
  \qquad
  S_n^{(k+1)}=\sum_{j=0}^{n}S_j^{(k)}.
  \label{eq:iterated-prefix-definition}
\end{equation}
Then \eqref{eq:inclusive-prefix} says
\begin{equation}
  S_{2q}^{(1)}=\eps_q,
  \qquad
  S_{2q+1}^{(1)}=0.
  \label{eq:first-inclusive-prefix}
\end{equation}

\begin{theorem}[Binomial convolution for discrete antiderivatives]
\label{thm:iterated-prefix-convolution}
For \(k\ge1\),
\begin{equation}
  \boxed{
  S_n^{(k)}=
  \sum_{j=0}^{n}
  \binom{n-j+k-1}{k-1}\eps_j.}
  \label{eq:iterated-prefix-convolution}
\end{equation}
Its ordinary generating function is
\begin{equation}
  \boxed{
  \sum_{n\ge0}S_n^{(k)}z^n
  =\frac{E(z)}{(1-z)^k}.}
  \label{eq:iterated-prefix-generating}
\end{equation}
\end{theorem}

\begin{proof}
One inclusive summation multiplies an ordinary generating function by
\((1-z)^{-1}\).  Iterate \(k\) times, then expand
\((1-z)^{-k}=\sum_{r\ge0}\binom{r+k-1}{k-1}z^r\) and take the Cauchy product.
\end{proof}

\subsection{Terminal zero runs at dyadic boundaries}

The quotient of a finite Prouhet polynomial by \((1-z)^k\) is still a
polynomial when \(k\le r\):
\begin{equation}
  \frac{P_r(z)}{(1-z)^k}
  =(1-z)^{r-k}
   \prod_{j=0}^{r-1}(1+z+\cdots+z^{2^j-1}).
  \label{eq:finite-prefix-polynomial}
\end{equation}
Its degree is \(2^r-k-1\).

\begin{theorem}[Dyadic terminal zero run]
\label{thm:terminal-zero-run}
If \(1\le k\le r\), then
\begin{equation}
  \boxed{
  S_{2^r-k+i}^{(k)}=0,
  \qquad 0\le i<k.}
  \label{eq:terminal-zero-run}
\end{equation}
The adjacent values are
\begin{equation}
  S_{2^r-k-1}^{(k)}=(-1)^{r-k},
  \qquad
  S_{2^r}^{(k)}=-1.
  \label{eq:terminal-zero-boundaries}
\end{equation}
\end{theorem}

\begin{proof}
Below degree \(2^r\), the coefficients of
\(E(z)/(1-z)^k\) agree with those of \(P_r(z)/(1-z)^k\).  The latter has
degree \(2^r-k-1\), proving the zero run.  Its leading coefficient is
\((-1)^{r-k}\), giving the preceding value.  The coefficient at \(2^r\)
comes from the first omitted factor \(1-z^{2^r}\), or directly from the prefix
recurrence.
\end{proof}

\subsection{Why this is the discrete spine of the Fabius theory}

The chain connecting Thue--Morse signs to the Fabius and Rvachev functions is
now algebraically visible:
\begin{equation}
 \boxed{\begin{gathered}
 \text{binary parity character}
 \Longrightarrow
 P_r(z)=\prod_{j<r}(1-z^{2^j})\\
 \Longrightarrow
 \dfrac{P_r(z)}{(1-z)^k}
 \Longrightarrow
 \text{iterated-prefix coefficient grids}
 \end{gathered}}
 \label{eq:fabius-discrete-chain}
\end{equation}
The zero of order \(k\) at \(z=1\) creates the exact terminal flatness in
\cref{thm:terminal-zero-run}; after the normalizations developed in the main
repository exposition, the resulting coefficient grids and polygonal
interpolants converge to the Fabius function.  On the Fourier side, replacing
\(z\) by a point on the unit circle turns the same finite product into a
finite product of sine factors, whose scaled limit is the sinc product
associated with the Rvachev up-function.

Thus the discrete-prefix and analytic-sinc appearances of Thue--Morse are not
separate phenomena.  They are the coefficient and Fourier-evaluation faces of
the single factorization \eqref{eq:finite-product}.  The quantitative
convergence theory belongs to the dedicated companion reports and Lean
modules; the purpose here is to expose the exact common algebraic spine.

\section[Fourier transforms]{Ordinary Fourier transforms and product-to-sum identities}
\label{sec:fourier}

\subsection{A finite Fourier product}

Putting \(z=\e^{\ii\theta}\) in \eqref{eq:finite-product} gives
\begin{equation}
  \sum_{n=0}^{2^m-1}\eps_n\e^{\ii n\theta}
  =\prod_{j=0}^{m-1}(1-\e^{\ii2^j\theta}).
  \label{eq:finite-Fourier-product-raw}
\end{equation}
Using \(1-\e^{\ii u}=-2\ii\e^{\ii u/2}\sin(u/2)\),
\begin{equation}
  \boxed{
  \sum_{n=0}^{2^m-1}\eps_n\e^{\ii n\theta}
  =(-2\ii)^m
   \e^{\ii(2^m-1)\theta/2}
   \prod_{j=0}^{m-1}\sin(2^{j-1}\theta).}
  \label{eq:finite-Fourier-sine}
\end{equation}
Hence
\begin{equation}
  \left|\sum_{n=0}^{2^m-1}\eps_n\e^{\ii n\theta}\right|
  =2^m\prod_{j=0}^{m-1}|\sin(2^{j-1}\theta)|.
  \label{eq:finite-Fourier-magnitude}
\end{equation}

\subsection{Sine and cosine products as Thue--Morse sums}

A rescaling of \eqref{eq:finite-Fourier-sine} gives exact finite
product-to-sum identities.

\begin{theorem}[Thue--Morse trigonometric sums]
\label{thm:trig-product-to-sum}
For every integer \(m\ge0\) and every complex \(z\),
\begin{equation}
  \boxed{
  \prod_{r=0}^{m}\sin\frac{z}{2^r}
  =\frac1{2^m}\sum_{n=0}^{2^m-1}\eps_n
   \sin\!\left(
     \frac{\pi m}{2}+\frac{(2n+1)z}{2^m}
   \right).}
  \label{eq:sine-product-to-sum}
\end{equation}
Also
\begin{equation}
  \boxed{
  \cos z\prod_{r=1}^{m}\sin\frac{z}{2^r}
  =\frac1{2^m}\sum_{n=0}^{2^m-1}\eps_n
   \cos\!\left(
     \frac{\pi m}{2}+\frac{(2n+1)z}{2^m}
   \right).}
  \label{eq:cosine-product-to-sum}
\end{equation}
The product from \(r=1\) to \(0\) is empty when \(m=0\).
\end{theorem}

\begin{proof}
Set \(\theta=2z/2^m\) in
\eqref{eq:finite-Fourier-sine}, multiply by
\(\ii^m\e^{\ii z/2^m}\), and simplify the phase.  The result is
\[
 \sum_{n<2^m}\eps_n
 \exp\!\left(\ii\left(
   \frac{\pi m}{2}+\frac{(2n+1)z}{2^m}
 \right)\right)
 =2^m\e^{\ii z}\prod_{r=1}^{m}\sin\frac z{2^r}.
\]
For real \(z\), take imaginary and real parts.  Both resulting identities
are identities between entire functions of \(z\), so the identity theorem
extends them to every complex \(z\).
\end{proof}

These identities make explicit the Fourier route from finite Thue--Morse
blocks to the sinc products in the Fabius/Rvachev analysis.

\subsection{The exact dyadic discrete Fourier transform}

Let
\begin{equation}
  \widehat\eps_m(k)=
  \sum_{n=0}^{2^m-1}\eps_n
  \exp\!\left(-\frac{2\pi\ii kn}{2^m}\right),
  \qquad 0\le k<2^m.
  \label{eq:DFT-definition}
\end{equation}

\begin{theorem}[Dyadic DFT]
\label{thm:dyadic-DFT}
For every \(k\),
\begin{equation}
  \widehat\eps_m(k)=
  \prod_{j=0}^{m-1}
  \left(1-\exp\!\left(-\frac{2\pi\ii k2^j}{2^m}\right)\right).
  \label{eq:DFT-product}
\end{equation}
If \(m\ge1\), this is zero for every even \(k\).  For odd \(k\),
\begin{equation}
  \boxed{
  \widehat\eps_m(k)=
  (2\ii)^m\e^{-\pi\ii k(1-2^{-m})}
  \prod_{r=1}^{m}\sin\frac{\pi k}{2^r}.}
  \label{eq:DFT-odd}
\end{equation}
Thus
\begin{equation}
  |\widehat\eps_m(k)|=
  2^m\prod_{r=1}^{m}\left|\sin\frac{\pi k}{2^r}\right|
  \qquad(k\text{ odd}).
  \label{eq:DFT-magnitude}
\end{equation}
\end{theorem}

\begin{proof}
Evaluate \eqref{eq:finite-product} at
\(z=\e^{-2\pi\ii k/2^m}\).  If \(k\) is even, the factor with
\(j=m-1\) is \(1-(-1)^k=0\).  For odd \(k\), apply
\(1-\e^{-\ii u}=2\ii\e^{-\ii u/2}\sin(u/2)\) factorwise and sum the phases.
\end{proof}

For \(m\ge1\), Fourier inversion gives a root-of-unity formula for every
individual term:
\begin{equation}
  \boxed{
  \eps_n=\frac1{2^m}
   \sum_{\substack{0\le k<2^m\\k\text{ odd}}}
   \widehat\eps_m(k)\e^{2\pi\ii kn/2^m},
  \qquad 0\le n<2^m.}
  \label{eq:DFT-inversion}
\end{equation}
This is intentionally extravagant as an evaluator, but exact and natural for
spectral calculations.

\section[Walsh, Riesz, and autocorrelation]{Walsh characters, Rademacher functions, Riesz products, and autocorrelation}
\label{sec:walsh-riesz}

\subsection{A one-point Walsh spectrum}

For \(m\)-bit vectors \(a,b\in\Ftwo^m\), write
\(\langle a,b\rangle=\sum_ja_jb_j\pmod2\), and let
\(\bm1=(1,\ldots,1)\).  For \(n<2^m\),
\begin{equation}
  \eps_n=(-1)^{\langle\bm1,\bm b(n)\rangle}.
  \label{eq:Walsh-character}
\end{equation}

\begin{theorem}[Walsh transform is a point mass]
\label{thm:Walsh-transform}
For every \(a\in\Ftwo^m\),
\begin{equation}
  \boxed{
  \sum_{n=0}^{2^m-1}\eps_n
  (-1)^{\langle a,\bm b(n)\rangle}
  =\begin{cases}
     2^m,&a=\bm1,\\
     0,&a\ne\bm1.
   \end{cases}}
  \label{eq:Walsh-transform}
\end{equation}
\end{theorem}

\begin{proof}
The summand is the character indexed by \(a+\bm1\).  A nontrivial character
sums to zero over the finite Boolean group, while the trivial character sums
to its order \(2^m\).
\end{proof}

Thus the broad ordinary Fourier spectrum is a basis-change phenomenon: the
same block is spectrally pure in the Walsh basis.

\subsection{Rademacher step functions}

For \(x\in[0,1)\), define the right-continuous Rademacher functions
\begin{equation}
  r_j(x)=(-1)^{\lfloor2^{j+1}x\rfloor}.
  \label{eq:Rademacher-definition}
\end{equation}
If \(0\le n<2^m\) and
\(x\in[n/2^m,(n+1)/2^m)\), then the first \(m\) binary digits of \(x\)
encode \(n\), so
\begin{equation}
  \boxed{
  \eps_n=\prod_{j=0}^{m-1}r_j(x).}
  \label{eq:Rademacher-product}
\end{equation}
The finite Thue--Morse block is therefore the sampling of one Walsh function
on the dyadic partition.

\subsection{Finite Riesz products}

Define
\begin{equation}
  \rho_m(x)=2^{-m}|P_m(\e^{2\pi\ii x})|^2.
  \label{eq:Riesz-density-definition}
\end{equation}

\begin{theorem}[Riesz-product density]
\label{thm:Riesz-product}
For every \(m\ge0\),
\begin{equation}
  \boxed{
  \rho_m(x)=\prod_{j=0}^{m-1}
  \bigl(1-\cos(2\pi2^jx)\bigr).}
  \label{eq:Riesz-product}
\end{equation}
Moreover,
\begin{equation}
  \int_0^1\rho_m(x)\dd x=1,
  \label{eq:Riesz-mass}
\end{equation}
so \(\rho_m(x)\dd x\) is a probability measure.
\end{theorem}

\begin{proof}
Use \(|1-\e^{\ii\theta}|^2=2(1-\cos\theta)\) in
\eqref{eq:finite-product}.  Parseval gives
\[
 \int_0^1\rho_m(x)\dd x
 =2^{-m}\sum_{n<2^m}|\eps_n|^2=1.
\]
\end{proof}

The finite measures converge weakly to the Thue--Morse spectral measure.  The
formal notation
\begin{equation}
  \dd\mu_{\mathrm{TM}}(x)
  =\prod_{j\ge0}(1-\cos(2\pi2^jx))\dd x
  \label{eq:infinite-Riesz-formal}
\end{equation}
means this weak limit, not pointwise multiplication of ordinary densities.
The limiting measure is purely singular continuous, a deeper theorem of
substitution spectral theory \cite{queffelec2010,baakeGrimm2008}.

\subsection{Finite and limiting autocorrelation}

For \(m,k\ge0\), define the dyadic average
\begin{equation}
  \eta_m(k)=\frac1{2^m}
   \sum_{n=0}^{2^m-1}\eps_n\eps_{n+k},
  \label{eq:finite-autocorrelation}
\end{equation}
where the second index is not reduced modulo \(2^m\).

\begin{proposition}[Finite autocorrelation recursion]
\label{prop:finite-autocorrelation}
For \(r\ge0\),
\begin{align}
  \eta_{m+1}(2r)&=\eta_m(r),
  \label{eq:eta-finite-even}\\
  \eta_{m+1}(2r+1)&=-\frac{\eta_m(r)+\eta_m(r+1)}2.
  \label{eq:eta-finite-odd}
\end{align}
\end{proposition}

\begin{proof}
Split the defining sum into even and odd \(n\).  For an even shift, both
halves reduce to \(\eps_j\eps_{j+r}\).  For an odd shift, they reduce to
\(-\eps_j\eps_{j+r}\) and \(-\eps_j\eps_{j+r+1}\), respectively.
\end{proof}

\begin{theorem}[Limiting autocorrelation]
\label{thm:limiting-autocorrelation}
For every fixed \(k\ge0\), the limit
\begin{equation}
  \eta(k)=\lim_{m\to\infty}\eta_m(k)
  \label{eq:eta-limit-definition}
\end{equation}
exists and is characterized by
\begin{align}
  \eta(0)&=1,
  \label{eq:eta-zero}\\
  \eta(2r)&=\eta(r),
  \label{eq:eta-even}\\
  \eta(2r+1)&=-\frac{\eta(r)+\eta(r+1)}2.
  \label{eq:eta-odd}
\end{align}
In particular,
\begin{equation}
  \eta(1)=-\frac13,
  \qquad
  \eta(2)=-\frac13,
  \qquad
  \eta(3)=\frac13,
  \qquad
  \eta(5)=0.
  \label{eq:eta-samples}
\end{equation}
\end{theorem}

\begin{proof}
Induct on \(k\).  The case \(k=0\) is constant.  For \(k=1\),
\eqref{eq:eta-finite-odd} gives the contraction
\[
 \eta_{m+1}(1)=-\frac{1+\eta_m(1)}2,
\]
whose fixed point is \(-1/3\).  If \(k=2r>1\), convergence reduces to the
smaller index \(r\).  If \(k=2r+1>1\), both \(r\) and \(r+1\) are smaller
than \(k\), so the odd recurrence and the induction hypothesis give
convergence.  Pass to the limit and evaluate recursively.
\end{proof}

The numbers \(\eta(k)\) are the Fourier coefficients of
\(\mu_{\mathrm{TM}}\).  Thus the substitution recursion, autocorrelation, and
Riesz product are exact combinatorial, coefficient, and measure-theoretic
forms of the same spectral object.

\section[Mellin--Dirichlet continuation]{Boundary flatness, shifted Dirichlet series, and Mellin continuation}
\label{sec:dirichlet-mellin}

This section develops an analytic consequence not organized systematically in
the four drafts.  The lacunary product is extraordinarily flat as its argument
approaches one from inside the unit disk.  That flatness removes every Mellin
singularity and gives an entire continuation of a natural family of
Thue--Morse Dirichlet series.

\subsection{Quadratic logarithmic decay near one}

For \(t>0\), set
\begin{equation}
  \mathcal E(t)=E(\e^{-t})
  =\prod_{j\ge0}(1-\e^{-2^jt}).
  \label{eq:E-caligraphic}
\end{equation}
Every factor lies in \((0,1)\), so \(\mathcal E(t)>0\).

\begin{theorem}[Boundary flatness]
\label{thm:boundary-flatness}
As \(t\downarrow0\),
\begin{equation}
  \boxed{
  \log\mathcal E(t)
  =-\frac{\log^2(1/t)}{2\log2}
   +\bigO(\log(1/t)).}
  \label{eq:boundary-log-asymptotic}
\end{equation}
Consequently, for every \(A>0\),
\begin{equation}
  \mathcal E(t)=\bigO_A(t^A)
  \qquad(t\downarrow0).
  \label{eq:superpolynomial-flatness}
\end{equation}
\end{theorem}

\begin{proof}
Assume \(0<t<1/2\), put \(L=\log(1/t)\), and let
\(m=\lfloor\log_2(1/t)\rfloor\).  Then
\(2^mt\in(1/2,1]\).  For \(0\le j<m\), one has \(0<2^jt\le1/2\), and the
elementary bounds
\begin{equation}
  \frac{x}{2}\le1-\e^{-x}\le x
  \qquad(0<x\le1)
  \label{eq:one-minus-exp-bounds}
\end{equation}
show that
\begin{equation}
  \log(1-\e^{-2^jt})=\log t+j\log2+\bigO(1),
  \label{eq:small-factor-log}
\end{equation}
with an absolute \(O(1)\).  Summing gives
\begin{equation}
 \sum_{j=0}^{m-1}\log(1-\e^{-2^jt})
 =m\log t+\frac{m(m-1)}2\log2+\bigO(m).
 \label{eq:small-factor-sum}
\end{equation}

For the tail, write \(2^mt=u\in(1/2,1]\).  Then
\begin{equation}
 \sum_{j=m}^{\infty}\left|\log(1-\e^{-2^jt})\right|
 =\sum_{r=0}^{\infty}\left|\log(1-\e^{-2^ru})\right|=\bigO(1)
 \label{eq:tail-uniform}
\end{equation}
uniformly in \(u\), because the terms decay doubly exponentially.  Finally
\(m=L/\log2+\bigO(1)\); substitution in
\eqref{eq:small-factor-sum} yields
\eqref{eq:boundary-log-asymptotic}.  The negative quadratic term dominates
\(-A\log(1/t)\) for every fixed \(A\), proving
\eqref{eq:superpolynomial-flatness}.
\end{proof}

\begin{remark}
Sharper asymptotic expansions contain a bounded log-periodic fluctuation in
\(\log_2(1/t)\).  The elementary estimate above is deliberately limited to the
part needed for Mellin continuation: super-polynomial flatness.
\end{remark}

\subsection{The shifted Dirichlet family}

For \(a>0\), define initially in the half-plane \(\Re s>0\)
\begin{equation}
  D(s,a)=\sum_{n=0}^{\infty}\frac{\eps_n}{(n+a)^s}.
  \label{eq:Dirichlet-definition}
\end{equation}
The principal real logarithm is used in \((n+a)^{-s}\).

\begin{proposition}[Conditional convergence]
\label{prop:Dirichlet-convergence}
The series \eqref{eq:Dirichlet-definition} converges locally uniformly in
\(\Re s>0\), and hence defines a holomorphic function there.
\end{proposition}

\begin{proof}
The partial sums \(S(N)=\sum_{n<N}\eps_n\) are bounded by one.  On a compact
set \(\Re s\ge\sigma>0\), the first difference of
\((n+a)^{-s}\) is uniformly \(\bigO(n^{-\sigma-1})\).  Summation by parts
therefore gives uniform convergence of the tails.
\end{proof}

Splitting even and odd indices gives an exact parameter equation.

\begin{theorem}[Dyadic Dirichlet equation]
\label{thm:Dirichlet-dyadic}
For \(a>0\) and \(\Re s>0\),
\begin{equation}
  \boxed{
  D(s,a)=2^{-s}\left(
    D\!\left(s,\frac a2\right)
    -D\!\left(s,\frac{a+1}{2}\right)
  \right).}
  \label{eq:Dirichlet-dyadic}
\end{equation}
\end{theorem}

\begin{proof}
Use \(\eps_{2n}=\eps_n\) and \(\eps_{2n+1}=-\eps_n\):
\begin{align*}
 D(s,a)
 &=\sum_{n\ge0}\frac{\eps_n}{(2n+a)^s}
   -\sum_{n\ge0}\frac{\eps_n}{(2n+1+a)^s}\\
 &=2^{-s}\left(
   D\!\left(s,\frac a2\right)
   -D\!\left(s,\frac{a+1}{2}\right)
 \right).
\end{align*}
\end{proof}

\subsection{Mellin representation and entire continuation}

\begin{theorem}[Mellin integral]
\label{thm:Mellin-integral}
For \(a>0\), initially when \(\Re s>1\),
\begin{equation}
  \boxed{
  D(s,a)=\frac1{\Gamma(s)}
  \int_0^{\infty}t^{s-1}\e^{-at}
  E(\e^{-t})\dd t.}
  \label{eq:Mellin-integral}
\end{equation}
The integral on the right is an entire function of \(s\).  Therefore
\eqref{eq:Mellin-integral} continues \(D(s,a)\) to an entire function of
\(s\in\C\).
\end{theorem}

\begin{proof}
For \(\Re s>1\), absolute convergence permits termwise use of
\[
 (n+a)^{-s}=\frac1{\Gamma(s)}
 \int_0^\infty t^{s-1}\e^{-(n+a)t}\dd t.
\]
Summing \(\eps_n\e^{-nt}\) gives \(E(\e^{-t})\).

At infinity, \(E(\e^{-t})=1+\bigO(\e^{-t})\), and the factor
\(\e^{-at}\) gives exponential decay.  At zero,
\cref{thm:boundary-flatness} gives decay faster than every positive power of
\(t\).  Thus the integral and all of its \(s\)-derivatives converge uniformly
on compact subsets of the entire \(s\)-plane.  Multiplication by the entire
function \(1/\Gamma(s)\) preserves entire-ness.  Equality on \(\Re s>1\),
together with \cref{prop:Dirichlet-convergence}, gives the continuation by the
identity theorem.
\end{proof}

\begin{corollary}[Zeros at nonpositive integers]
\label{cor:Dirichlet-trivial-zeros}
For every \(a>0\) and every integer \(r\ge0\),
\begin{equation}
  \boxed{D(-r,a)=0.}
  \label{eq:Dirichlet-trivial-zeros}
\end{equation}
\end{corollary}

\begin{proof}
The Mellin integral in \eqref{eq:Mellin-integral} is finite at \(s=-r\),
while \(1/\Gamma(s)\) has a zero there.
\end{proof}

These zeros are the analytic continuation of the finite Prouhet phenomenon:
negative integer values formally ask for signed polynomial moments, and the
flat Mellin kernel makes all of them vanish after continuation.

\subsection{The harmonic kernel}

At \(s=1\), the Mellin integral can be moved to the unit interval.

\begin{corollary}[Dirichlet integral]
\label{cor:harmonic-integral}
For \(a>0\),
\begin{align}
  H(a):=D(1,a)
  &=\sum_{n\ge0}\frac{\eps_n}{n+a}
  \label{eq:H-definition}\\
  &=\boxed{
  \int_0^1x^{a-1}E(x)\dd x
  =\int_0^1x^{a-1}
    \prod_{j\ge0}(1-x^{2^j})\dd x}>0.
  \label{eq:H-integral}
\end{align}
It satisfies
\begin{equation}
  \boxed{
  H(a)=\frac12H\!\left(\frac a2\right)
       -\frac12H\!\left(\frac{a+1}{2}\right).}
  \label{eq:H-dyadic}
\end{equation}
\end{corollary}

\begin{proof}
Set \(s=1\) in \eqref{eq:Mellin-integral} and substitute \(x=\e^{-t}\).
The integrand is positive on \((0,1)\).  Equation \eqref{eq:H-dyadic} is
\eqref{eq:Dirichlet-dyadic} at \(s=1\).
\end{proof}

\section[Thue--Morse exponent products]{Infinite products with Thue--Morse exponents}
\label{sec:exponent-products}

\subsection{Pair reduction}

Suppose \(R(n)\ne0\) and the products converge.  The even--odd recurrence
gives the general reduction
\begin{equation}
  \boxed{
  \prod_{n\ge0}R(n)^{\eps_n}
  =\prod_{q\ge0}
    \left(\frac{R(2q)}{R(2q+1)}\right)^{\eps_q}.}
  \label{eq:pair-reduction}
\end{equation}
A factor merely equal to \(1+\bigO(1/n)\) is often converted into a paired
factor \(1+\bigO(1/n^2)\), explaining the good convergence of balanced
rational products.

\subsection{A two-parameter master product}

For \(a,b>0\), define
\begin{equation}
  G(a,b)=\prod_{n=0}^{\infty}
  \left(\frac{n+a}{n+b}\right)^{\eps_n}.
  \label{eq:G-definition}
\end{equation}
The logarithm converges by summation by parts: the Thue--Morse partial sums are
bounded and the first difference of
\(\log((n+a)/(n+b))\) is \(\bigO(n^{-2})\).

\begin{theorem}[Cocycle and dyadic functional equation]
\label{thm:G-functional}
For positive parameters,
\begin{align}
  G(a,a)&=1,
  &G(a,b)G(b,c)&=G(a,c),
  &G(b,a)&=G(a,b)^{-1},
  \label{eq:G-cocycle}\\
  G(a,b)&=\boxed{
  \frac{G(a/2,b/2)}
       {G((a+1)/2,(b+1)/2)}}.
  \label{eq:G-functional}
\end{align}
\end{theorem}

\begin{proof}
The cocycle identities are termwise.  Pair indices \(2n\) and \(2n+1\):
\begin{align*}
 &\left(\frac{2n+a}{2n+b}\right)^{\eps_n}
 \left(\frac{2n+1+a}{2n+1+b}\right)^{-\eps_n}\\
 &\qquad=
 \left[
  \frac{n+a/2}{n+b/2}
  \frac{n+(b+1)/2}{n+(a+1)/2}
 \right]^{\eps_n}.
\end{align*}
Multiply over \(n\).
\end{proof}

The Dirichlet continuation identifies this product as a zeta-regularized
derivative difference without any extra regularization being required.

\begin{corollary}[Dirichlet derivative bridge]
\label{cor:G-Dirichlet}
For \(a,b>0\),
\begin{equation}
  \boxed{
  \log G(a,b)=D'(0,b)-D'(0,a).}
  \label{eq:G-Dirichlet}
\end{equation}
\end{corollary}

\begin{proof}
The difference \(D(s,a)-D(s,b)\) can be differentiated termwise at
\(s=0\), because the resulting log-ratio series converges by summation by
parts.  Its derivative is
\[
 \sum_{n\ge0}\eps_n
 \log\frac{n+b}{n+a}=-\log G(a,b).
\]
Rearrange.
\end{proof}

\subsection{The Woods--Robbins product}

\begin{theorem}[Woods--Robbins]
\label{thm:Woods-Robbins}
\begin{equation}
  \boxed{
  \prod_{n=0}^{\infty}
  \left(\frac{2n+1}{2n+2}\right)^{\eps_n}
  =\frac1{\sqrt2}.}
  \label{eq:Woods-Robbins}
\end{equation}
\end{theorem}

\begin{proof}
Let
\begin{equation}
 A=\prod_{n\ge0}
   \left(\frac{2n+1}{2n+2}\right)^{\eps_n},
 \qquad
 B=\prod_{n\ge1}
   \left(\frac{2n}{2n+1}\right)^{\eps_n}.
 \label{eq:A-B-products}
\end{equation}
Then
\begin{equation}
 AB=\frac12\prod_{n\ge1}
       \left(\frac n{n+1}\right)^{\eps_n}.
 \label{eq:AB-first}
\end{equation}
Split the product on the right into odd and even \(n\).  The odd part is
\(A^{-1}\) because \(\eps_{2j+1}=-\eps_j\), and the even part is \(B\)
because \(\eps_{2j}=\eps_j\).  Hence
\begin{equation}
  AB=\frac12A^{-1}B.
\end{equation}
Cancel \(B\), obtaining \(A^2=1/2\); positivity gives the stated root.
\end{proof}

A second classical evaluation is
\begin{equation}
  \boxed{
  \prod_{n=0}^{\infty}
  \left(\frac{4n+1}{4n+3}\right)^{\eps_n}=\frac12.}
  \label{eq:quarter-product}
\end{equation}
It is the specialization \(G(1/4,3/4)\) and follows from the broader
Thue--Morse gamma-product calculus of Allouche, Riasat, and Shallit
\cite{alloucheRiasatShallit2019}.

\subsection{Three dyadic block-product limits}

Dyadic truncation respects exact block recursion.  The following limits are
convenient fingerprints of the same product theory:
\begin{align}
  \lim_{m\to\infty}
  \prod_{k=0}^{2^m-1}\left(k+\frac12\right)^{\eps_k}
  &=\frac12,
  \label{eq:block-product-half}\\
  \lim_{m\to\infty}
  \prod_{k=0}^{2^m-1}(k+1)^{\eps_k}
  &=\frac1{\sqrt2},
  \label{eq:block-product-one}\\
  \lim_{m\to\infty}
  \prod_{k=0}^{2^m-1}(k+1)^{(-1)^k\eps_k}
  &=\frac1{2\sqrt2}.
  \label{eq:block-product-alternating}
\end{align}
Pairing adjacent factors reduces the first two to
\(G(1/4,3/4)\) and \(G(1/2,1)\), respectively; the third factors into their
product.  The evaluations are thus \eqref{eq:quarter-product} and
\eqref{eq:Woods-Robbins}.

\section[Base-\texorpdfstring{\(q\)}{q} generalization]{Base-\texorpdfstring{\(q\)}{q} digit characters and generalized Prouhet identities}
\label{sec:base-q}

Let \(q\ge2\), let \(s_q(n)\) be the sum of the base-\(q\) digits of \(n\),
and let \(\zeta\ne1\) satisfy \(\zeta^q=1\).  Define
\begin{equation}
  u_n=\zeta^{s_q(n)}.
  \label{eq:generalized-digit-character}
\end{equation}

\begin{theorem}[Generalized digit product]
\label{thm:base-q-product}
For \(0\le r<q\),
\begin{equation}
  u_{qn+r}=\zeta^ru_n.
  \label{eq:base-q-recursion}
\end{equation}
For every \(m\ge0\),
\begin{equation}
  \boxed{
  \sum_{n=0}^{q^m-1}u_nz^n
  =\prod_{j=0}^{m-1}
  \left(\sum_{r=0}^{q-1}\zeta^rz^{rq^j}\right).}
  \label{eq:base-q-product}
\end{equation}
\end{theorem}

\begin{proof}
Appending the base-\(q\) digit \(r\) increases the digit sum by \(r\).  In the
product, choose one digit independently at each position; the exponent is the
corresponding base-\(q\) integer and the coefficient is \(\zeta\) raised to
its digit sum.
\end{proof}

Each factor vanishes at \(z=1\), since
\(1+\zeta+\cdots+\zeta^{q-1}=0\).

\begin{theorem}[Generalized Prouhet cancellation and sharp moment]
\label{thm:base-q-Prouhet}
If \(p\) is a polynomial of degree less than \(m\), then
\begin{equation}
  \boxed{
  \sum_{n=0}^{q^m-1}\zeta^{s_q(n)}p(n)=0.}
  \label{eq:base-q-Prouhet}
\end{equation}
At the first surviving degree,
\begin{equation}
  \boxed{
  \sum_{n=0}^{q^m-1}\zeta^{s_q(n)}n^m
  =m!\frac{q^{m(m+1)/2}}{(\zeta-1)^m}.}
  \label{eq:base-q-sharp-moment}
\end{equation}
\end{theorem}

\begin{proof}
Set \(z=\e^t\) in \eqref{eq:base-q-product}.  Every factor has a zero at
\(t=0\), proving cancellation below degree \(m\).  Its derivative at zero is
\begin{equation}
 q^j\sum_{r=0}^{q-1}r\zeta^r
 =q^j\frac{q}{\zeta-1}
 =\frac{q^{j+1}}{\zeta-1},
 \label{eq:root-unity-derivative}
\end{equation}
where the finite geometric sum was differentiated.  Multiplying the \(m\)
linear terms yields
\(q^{1+2+\cdots+m}/(\zeta-1)^m\); coefficient extraction contributes \(m!\).
\end{proof}

For \(q=2\) and \(\zeta=-1\), the sequence is \(u_n=\eps_n\) and
\eqref{eq:base-q-sharp-moment} becomes
\((-1)^mm!2^{\binom m2}\), exactly
\eqref{eq:first-surviving-moment}.  The binary theory is therefore the first
member of a root-of-unity digit-character family, while the particularly rich
valuation, carry, and Pascal-row formulas remain genuinely base-two features.

\section[Formula selection and Lean roadmap]{Formula selection, computational checks, and a Lean roadmap}
\label{sec:formalization}

\subsection{Which representation is best?}

The formulas are equivalent mathematically but not proof-theoretically or
computationally.  The following table records the most economical interface
for common tasks.

\renewcommand{\arraystretch}{1.16}
\begin{longtable}{@{}>{\raggedright\arraybackslash}p{0.22\textwidth}>{\raggedright\arraybackslash}p{0.32\textwidth}>{\raggedright\arraybackslash}p{0.38\textwidth}@{}}
\toprule
Goal & Recommended representation & Reason \\
\midrule
\endfirsthead
\toprule
Goal & Recommended representation & Reason \\
\midrule
\endhead
Executable evaluation & Recurrence or machine popcount & Linear in the bit length; no large intermediate integers. \\
Dyadic reindexing & Block concatenation \eqref{eq:block-sign} & Quotient and remainder separate without carries. \\
Bitwise algebra & XOR character \eqref{eq:xor-character} & Makes \(\eps\) a genuine group character. \\
Ordinary addition & Carry cocycle \eqref{eq:carry-sign} & Isolates the exact obstruction to XOR multiplicativity. \\
Number-theoretic disguise & Central binomial valuation \eqref{eq:central-thue-morse} & One classical integer stores the full binary weight. \\
Pascal-row interpretation & Glaisher count or modulo-three residue & Uses the odd support of one row; the residue form avoids logarithms. \\
Higher-dimensional Pascal & Multinomial count \eqref{eq:multinomial-support} & Works with any number of variables. \\
Formal power series & Finite/infinite products & Coefficients, Mahler equations, and zeros become immediate. \\
Nonlocal coefficient formulas & Euler/Bell/determinant forms & Eliminate explicit digit notation and expose ruler convolutions. \\
Polynomial cancellation & Mixed finite differences & Proves all Prouhet moments in one operator identity. \\
Arbitrary Pascal row & Sparse signed product \eqref{eq:signed-Pascal-support} & Restricts finite differences to the odd row support. \\
Finite harmonic analysis & Walsh character & The entire block is one pure frequency. \\
Ordinary spectral analysis & DFT sine product and Riesz density & Exposes roots of unity, amplitudes, and correlations. \\
Analytic continuation & Mellin integral \eqref{eq:Mellin-integral} & Boundary flatness removes every Mellin singularity. \\
Infinite products & Master product \(G(a,b)\) & Pair reduction and cocycle laws organize evaluations. \\
Fabius approximation & Iterated-prefix generating function & Coefficient grids arise from \(E(z)/(1-z)^k\). \\
\bottomrule
\end{longtable}
\renewcommand{\arraystretch}{1}

\subsection{A compact verification script}

The following pseudocode is not a proof, but it is useful for regression tests
when translating the atlas into Lean, a CAS, or a library implementation.

\begin{lstlisting}[style=pseudo]
for n from 0 to N:
    s := popcount(n)
    eps := 1 - 2*(s mod 2)

    assert v2(factorial(n)) == n - s
    assert v2(binomial(2*n,n)) == s

    oddPascal := count k in [0,n] with binomial(n,k) odd
    assert oddPascal == 2^s
    assert (oddPascal mod 3) - 1 == (s mod 2)

    sparseParity := sum binomial(n,2^j) over all 2^j <= n
    assert (-1)^sparseParity == eps

for m from 0 to M:
    P := sum_{n=0}^{2^m-1} eps(n)*z^n
    assert P == product_{j=0}^{m-1} (1-z^(2^j))

    for r from 0 to m-1:
        assert sum_{n=0}^{2^m-1} eps(n)*n^r == 0

    assert sum_{n=0}^{2^m-1} eps(n)*n^m
           == (-1)^m * factorial(m) * 2^(m*(m-1)/2)
\end{lstlisting}

\subsection{A dependency-aware formalization order}

A conservative Lean implementation should formalize the finite algebraic
spine before introducing analysis.

\begin{enumerate}
\item \textbf{Binary layer.}  Prove recurrence, block concatenation, XOR,
      reflection, and the two-kernel from a chosen binary-weight primitive.
\item \textbf{Valuation layer.}  Establish the base-two Legendre formula,
      central-binomial valuation, Catalan valuation, and Kummer carry identity.
\item \textbf{Pascal/Frobenius layer.}  Formalize Lucas support, Glaisher count,
      sparse probes, the modulo-three consequence, multinomial support, and
      Boolean M\"obius inversion.
\item \textbf{Polynomial layer.}  Define \(P_m\), prove the finite product,
      reciprocity, cyclotomic multiplicities, and signed Pascal-support
      products over a suitable coefficient ring.
\item \textbf{Difference and moment layer.}  Derive mixed finite differences,
      Prouhet cancellation, sharp moments, composition formulas, and sparse
      moments.  Most statements can be done over \(\Q\) before any topology is
      imported.
\item \textbf{Formal-series layer.}  Add the Mahler equations, Euler transform,
      ruler convolution, Bell/determinant forms, and the Artin--Schreier
      equation over \(\Ftwo[[z]]\).
\item \textbf{Analytic layer.}  Only then introduce compact-uniform convergence,
      Fourier evaluation, Riesz measures, Mellin transforms, and infinite
      products.
\end{enumerate}

Possible module boundaries are
\begin{center}
\small
\begin{tabularx}{0.94\textwidth}{@{}>{\raggedright\arraybackslash}p{0.42\textwidth}>{\raggedright\arraybackslash}X@{}}
\toprule
Proposed module & Main declarations \\
\midrule
\path{ThueMorseValuation.lean} & Factorial, central-binomial, Catalan, successor, and carry formulae. \\
\path{ThueMorsePascal.lean} & Lucas submasks, Glaisher count, sparse probes, modulo-three formula. \\
\path{ThueMorseMultinomial.lean} & Frobenius factorization and \(O_r(n)=r^{\wt(n)}\). \\
\path{ThueMorseBooleanMobius.lean} & Submask incidence algebra and inverse parity matrix. \\
\path{ThueMorseGeneratingProduct.lean} & \(P_m\), cyclotomic multiplicities, Mahler equations. \\
\path{ThueMorseProuhetMoments.lean} & Difference operators, all moments, centered parity. \\
\path{ThueMorseSparseProuhet.lean} & Arbitrary Pascal-support products and sparse moments. \\
\path{ThueMorseFormalSeriesF2.lean} & Artin--Schreier equation and explicit solution. \\
\path{ThueMorseFourier.lean} & Finite DFT, Walsh transform, product-to-sum identities. \\
\path{ThueMorseDirichlet.lean} & Boundary flatness, Mellin integral, entire continuation. \\
\bottomrule
\end{tabularx}
\end{center}

The last module is analytically the heaviest.  The first seven are finite or
formal and can deliver useful lemmas to the existing Fabius development
without waiting for measure theory or complex Mellin infrastructure.

\subsection{Proof-dependency graph}

The central dependencies can be summarized as
\begin{equation}
\begin{aligned}
 \text{binary digits}
 &\Longrightarrow
 \begin{cases}
   \text{recursion and XOR},\\
   \text{Lucas/Frobenius support},\\
   \text{finite product }P_m,
 \end{cases}\\[0.4em]
 P_m
 &\Longrightarrow
 \begin{cases}
   \text{Mahler and Euler equations},\\
   \text{Fourier/Walsh products},\\
   \text{order-}m\text{ zero at }1,
 \end{cases}\\[0.4em]
 \text{order-}m\text{ zero}
 &\Longrightarrow
 \begin{cases}
   \text{Prouhet cancellation},\\
   \text{all finite moments},\\
   \text{iterated-prefix flatness},\\
   \text{Riesz-product normalization}.
 \end{cases}
\end{aligned}
\label{eq:dependency-graph}
\end{equation}
Legendre/Kummer form a parallel arithmetic branch; Lucas and the signed
support product join that branch to the finite-difference branch at sparse
Prouhet identities.  Boundary flatness of the infinite product then opens the
Mellin/Dirichlet branch.

\section[Worked examples]{Worked examples and consistency checks}
\label{sec:worked-examples}

\subsection{One index through many representations}

Take \(n=13=(1101)_2\).  The occupied positions are \(0,2,3\), so
\begin{equation}
  \wt(13)=3,
  \qquad
  \tau(13)=1,
  \qquad
  \eps_{13}=-1.
  \label{eq:example13-basic}
\end{equation}
The floor, factorial, and central-binomial routes give
\begin{align}
 \sum_{j\ge1}\left\lfloor\frac{13}{2^j}\right\rfloor
 &=6+3+1=10,
 \label{eq:example13-floor}\\
 (-1)^{13-\vtwo(13!)}&=(-1)^3=-1,
 \label{eq:example13-factorial}\\
 \vtwo\binom{26}{13}&=3.
 \label{eq:example13-central}
\end{align}
The distinguished Pascal entries are
\begin{equation}
 \binom{13}{1}=13,
 \quad
 \binom{13}{2}=78,
 \quad
 \binom{13}{4}=715,
 \quad
 \binom{13}{8}=1287.
 \label{eq:example13-probes}
\end{equation}
Their parities are \(1,0,1,1\), exactly the four low bits of \(13\), and their
sum \(2093\) is odd.  Row \(13\) has
\begin{equation}
  \nu(13)=2^3=8
  \label{eq:example13-odd-count}
\end{equation}
odd entries, so \(\nu(13)\bmod3=2\) and the logarithm-free formula gives
\(\tau(13)=2-1=1\).  In three variables, exactly
\begin{equation}
  O_3(13)=3^3=27
  \label{eq:example13-multinomial}
\end{equation}
coefficients of \((x+y+z)^{13}\) are odd.

The submasks are
\begin{equation}
  0,1,4,5,8,9,12,13,
  \label{eq:example13-submasks}
\end{equation}
and their signed support polynomial is
\begin{equation}
  R_{13}(z)=(1-z)(1-z^4)(1-z^8).
  \label{eq:example13-support-product}
\end{equation}
Its derivatives of orders \(0,1,2\) at one vanish, and
\begin{equation}
  \sum_{k\submask13}\eps_kk^3=-3!\,2^{0+2+3}=-192.
  \label{eq:example13-sparse-moment}
\end{equation}
All these evaluations recover the same three occupied bits through different
algebraic interfaces.

\subsection{An all-order dyadic moment}

For \(m=3\) and \(r=5\), direct summation gives
\begin{equation}
  M_{3,5}=\sum_{n=0}^{7}\eps_nn^5=-6720.
  \label{eq:example-moment-direct}
\end{equation}
The positive compositions of \(5\) into three parts are
\[
 (1,1,3),(1,2,2),(1,3,1),(2,1,2),(2,2,1),(3,1,1).
\]
Formula \eqref{eq:moment-positive-compositions} gives
\begin{equation}
  M_{3,5}
  =-5!\sum_{\substack{p_0+p_1+p_2=5\\p_j\ge1}}
     \frac{2^{p_1+2p_2}}{p_0!p_1!p_2!}
  =-6720.
  \label{eq:example-moment-composition}
\end{equation}
No recursion through lower moments is needed.

\subsection{A Fourier sample}

For \(m=4\), every even DFT frequency vanishes.  At \(k=1\),
\begin{equation}
  \widehat\eps_4(1)
  =(2\ii)^4\e^{-15\pi\ii/16}
  \sin\frac\pi2\sin\frac\pi4
  \sin\frac\pi8\sin\frac\pi{16}.
  \label{eq:example-Fourier}
\end{equation}
The factorization exposes both the exact phase and the hierarchy of four
dyadic scales, which direct summation of sixteen complex terms obscures.

\subsection{Autocorrelation from the recursion}

Starting with \(\eta(0)=1\), the fixed-point equation at one gives
\(\eta(1)=-1/3\).  Then
\begin{align*}
 \eta(2)&=\eta(1)=-\frac13,\\
 \eta(3)&=-\frac{\eta(1)+\eta(2)}2=\frac13,\\
 \eta(5)&=-\frac{\eta(2)+\eta(3)}2=0.
\end{align*}
These are simultaneously combinatorial averages and Fourier coefficients of
the singular continuous Thue--Morse measure.

\section[Conclusion]{Concluding synthesis}
\label{sec:conclusion}

The formula in Section~11.8 of the Fabius--Rvachev exposition is not an
isolated trick.  It lies at the intersection of four structures:
\begin{equation}
 \boxed{\begin{gathered}
 \text{binary digits}
 \longleftrightarrow
 \text{Frobenius/Pascal support}\\
 \longleftrightarrow
 \text{Boolean-lattice characters}
 \longleftrightarrow
 \text{dyadic finite differences}
 \end{gathered}}
 \label{eq:four-way-dictionary}
\end{equation}
The signed Thue--Morse sequence is the common coefficient system of this
translation dictionary.

The arithmetic formulas hide digit parity in factorial valuations, central
binomial coefficients, Catalan numbers, and carry counts.  The Pascal formulas
recover it either globally, from the number of odd row entries, or sparsely,
from the power-of-two columns.  The modulo-three identity shows that the exact
base-two logarithm is optional: the exponent parity can also be read from the
cycle of powers of two in a finite residue field.  Multinomial support turns the
same mechanism into a whole \(r^{\wt(n)}\) family.

The product
\(P_m(z)=\prod_{j<m}(1-z^{2^j})\) is the central unifying engine.  At
\(z=1\) it has an exact high-order zero and produces Prouhet cancellation.  At
\(z=\e^t\) it produces every finite moment.  At roots of unity it produces the
dyadic DFT.  On the Boolean cube it is a pure Walsh character.  After taking
absolute squares it produces Riesz densities.  After division by
\((1-z)^k\) it produces iterated-prefix grids and their terminal flatness.
After the limit \(m\to\infty\), it becomes a Mahler function, an Euler
transform of the ruler sequence, and the Mellin kernel of an entire Dirichlet
family.

The added Mellin theory sharpens the analytic picture: the infinite product is
not merely small near one, but flat beyond every algebraic order, with a
quadratic leading term in \(\log(1/t)\).  That single estimate explains why
all potential Mellin poles disappear and why the shifted Thue--Morse Dirichlet
series has zeros at every nonpositive integer.  Its derivative at zero then
recovers the logarithms of the gamma-like products generated by Thue--Morse
exponents.

For the Fabius project, the most useful lesson is proof-interface diversity.
No one formula dominates every task.  Block multiplicativity is ideal for
dyadic reindexing; the carry cocycle for ordinary sums; Pascal support for
characteristic-two arguments; finite differences for cancellation; Bell
polynomials for all moments; Walsh characters for finite harmonic analysis;
and the Mellin product for analytic continuation.  Formalizing several of
these interfaces would not duplicate mathematics: it would let later Fabius
proofs choose the representation adapted to the operation being performed.

\appendix

\section[First 64 terms]{The first 64 terms}
\label{app:first-terms}

The rows display \(\tau(n)\); replacing \(0\) by \(+1\) and \(1\) by
\(-1\) gives \(\eps_n\).

\begin{center}
\small
\begin{longtable}{@{}r|*{8}{c}@{}}
\toprule
Range & \multicolumn{8}{c}{successive values}\\
\midrule
\endhead
0--7   & 0&1&1&0&1&0&0&1\\
8--15  & 1&0&0&1&0&1&1&0\\
16--23 & 1&0&0&1&0&1&1&0\\
24--31 & 0&1&1&0&1&0&0&1\\
32--39 & 1&0&0&1&0&1&1&0\\
40--47 & 0&1&1&0&1&0&0&1\\
48--55 & 0&1&1&0&1&0&0&1\\
56--63 & 1&0&0&1&0&1&1&0\\
\bottomrule
\end{longtable}
\end{center}

Every row of length eight is a copy or complement according to the sign of its
block index, as predicted by \eqref{eq:block-sign}.

\section[Identity dictionary]{A compact identity dictionary}
\label{app:identity-dictionary}

For quick reference, the principal exact translations are
\begin{align*}
 \tau(n)
 &\longleftrightarrow \eps_n=1-2\tau(n),\\
 \eps_n
 &\longleftrightarrow (-1)^{\wt(n)},\\
 \wt(n)
 &\longleftrightarrow n-\vtwo(n!),\\
 \wt(n)
 &\longleftrightarrow \vtwo\binom{2n}{n},\\
 b_j(n)
 &\longleftrightarrow \binom{n}{2^j}\bmod2,\\
 2^{\wt(n)}
 &\longleftrightarrow \#\{k:\binom nk\text{ odd}\},\\
 \eps_n
 &\longleftrightarrow [z^n]\prod_{j\ge0}(1-z^{2^j}),\\
 \sum_{n<2^m}\eps_nf(x+n)
 &\longleftrightarrow (-1)^m\Delta_1\Delta_2\cdots
   \Delta_{2^{m-1}}f(x),\\
 P_m(\e^{-2\pi\ii k/2^m})
 &\longleftrightarrow \widehat\eps_m(k),\\
 E(\e^{-t})
 &\longleftrightarrow \Gamma(s)D(s,a)
  \text{ under Mellin transformation}.
\end{align*}
Every arrow denotes an exact identity or invertible conversion, not an
asymptotic analogy.

\section[Draft concordance]{Draft-to-section concordance}
\label{app:draft-concordance}

The four source drafts overlap substantially.  The following table records
the main distinctive material retained from each, together with additions made
in the unified treatment.

\renewcommand{\arraystretch}{1.16}
\begin{longtable}{@{}p{0.20\textwidth}p{0.72\textwidth}@{}}
\toprule
Source & Distinctive contributions represented in this atlas \\
\midrule
\endfirsthead
\toprule
Source & Distinctive contributions represented in this atlas \\
\midrule
\endhead
Draft 1 & Carry and XOR viewpoints; Boolean M\"obius inversion and inverse Pascal parity matrix; ruler Euler transform; partition, Bell, and Hessenberg determinant formulae; finite-difference positivity and analytic transforms. \\
Draft 2 & Sparse power-of-two Pascal columns; all-order dyadic moments; automaton/tensor/Walsh forms; exact DFT and Riesz products; master infinite products; base-\(q\) root-of-unity extension. \\
Draft 3 & Multinomial support and exact \(r\)-nomial logarithms; sparse Prouhet theorem on arbitrary Pascal rows; autocorrelation recurrences; Artin--Schreier series over \(\mathbb F_2\); Catalan and fractional-part formulae. \\
Draft 4 & Logarithm-free modulo-three Pascal formula; sinc-product interval signs; evil/odious enumerators; integer algebraic lift; trigonometric product-to-sum identities; explicit iterated-prefix/Fabius bridge. \\
Unified additions & Dyadic weight-distribution chapter; centered-moment parity selection; all sparse Pascal-support moments; elementary boundary asymptotic for \(E(\e^{-t})\); entire shifted Dirichlet family, its nonpositive-integer zeros, and its derivative bridge to exponent products; sharp first moment for the general base-\(q\) character. \\
\bottomrule
\end{longtable}
\renewcommand{\arraystretch}{1}

\begin{thebibliography}{99}

\bibitem{proveit}
V.~Reshetnikov,
\newblock \emph{ProveIt: Analysis/FabiusFunction},
\newblock GitHub repository, source snapshot
\nolinkurl{29e5861c6109207739c41754e1e8bcfce32e84f1}, 2026.

\bibitem{prouhet1851}
E.~Prouhet,
\newblock M\'emoire sur quelques relations entre les puissances des nombres,
\newblock \emph{Comptes rendus de l'Acad\'emie des Sciences de Paris}
\textbf{33} (1851), 225.

\bibitem{thue1906}
A.~Thue,
\newblock \emph{\"Uber unendliche Zeichenreihen},
\newblock Norske Videnskabers Selskabs Skrifter, I. Matematisk-naturvidenskabelig
Klasse, no.~7 (1906), 1--22.

\bibitem{morse1921}
M.~Morse,
\newblock Recurrent geodesics on a surface of negative curvature,
\newblock \emph{Transactions of the American Mathematical Society}
\textbf{22} (1921), 84--100.

\bibitem{kummer1852}
E.~E.~Kummer,
\newblock \"Uber die Erg\"anzungss\"atze zu den allgemeinen
Reciprocit\"atsgesetzen,
\newblock \emph{Journal f\"ur die reine und angewandte Mathematik}
\textbf{44} (1852), 93--146.

\bibitem{fine1947}
N.~J.~Fine,
\newblock Binomial coefficients modulo a prime,
\newblock \emph{American Mathematical Monthly} \textbf{54} (1947), 589--592.

\bibitem{mestrovi\'c2014}
R.~Me\v{s}trovi\'c,
\newblock Lucas' theorem: its generalizations, extensions and applications
(1878--2014),
\newblock arXiv:\nolinkurl{1409.3820}, 2014.

\bibitem{mahler1929}
K.~Mahler,
\newblock Arithmetische Eigenschaften der L\"osungen einer Klasse von
Funktionalgleichungen,
\newblock \emph{Mathematische Annalen} \textbf{101} (1929), 342--366.

\bibitem{woods1978}
D.~R.~Woods,
\newblock Elementary problem E2692,
\newblock \emph{American Mathematical Monthly} \textbf{85} (1978), 48.

\bibitem{robbins1979}
D.~Robbins,
\newblock Solution to elementary problem E2692,
\newblock \emph{American Mathematical Monthly} \textbf{86} (1979), 394--395.

\bibitem{alloucheShallit1999}
J.-P.~Allouche and J.~Shallit,
\newblock The ubiquitous Prouhet--Thue--Morse sequence,
\newblock in C.~Ding, T.~Helleseth, and H.~Niederreiter (eds.),
\emph{Sequences and Their Applications}, Springer, London, 1999, 1--16.

\bibitem{alloucheShallit2003}
J.-P.~Allouche and J.~Shallit,
\newblock \emph{Automatic Sequences: Theory, Applications, Generalizations},
\newblock Cambridge University Press, Cambridge, 2003.

\bibitem{christol1980}
G.~Christol, T.~Kamae, M.~Mend\`es France, and G.~Rauzy,
\newblock Suites alg\'ebriques, automates et substitutions,
\newblock \emph{Bulletin de la Soci\'et\'e Math\'ematique de France}
\textbf{108} (1980), 401--419.

\bibitem{queffelec2010}
M.~Queff\'elec,
\newblock \emph{Substitution Dynamical Systems---Spectral Analysis},
2nd ed., Lecture Notes in Mathematics 1294,
\newblock Springer, Berlin, 2010.

\bibitem{baakeGrimm2008}
M.~Baake and U.~Grimm,
\newblock The singular continuous diffraction measure of the Thue--Morse chain,
\newblock \emph{Journal of Physics A: Mathematical and Theoretical}
\textbf{41} (2008), 422001.

\bibitem{alloucheRiasatShallit2019}
J.-P.~Allouche, S.~Riasat, and J.~Shallit,
\newblock More infinite products: Thue--Morse and the gamma function,
\newblock \emph{Ramanujan Journal} \textbf{49} (2019), 115--128;
\newblock arXiv:\nolinkurl{1709.03398}.

\end{thebibliography}

\end{document}
